\documentclass[../../book-main.tex]{subfiles}
\usepackage[UTF8]{ctex}
\graphicspath{{\subfix{../..}}}

\begin{document}

\chapter{通过压缩学习低维分布}
\label{ch:compression}\label{ch:general-distribution}

\begin{quote}
\hfill    ``{\em 我们通过压缩来学习，我们通过学习来压缩}.''\\
$~$ \hfill --- 高维数据分析，2022   
\end{quote}
\vspace{5mm}

在第 \ref{ch:linear-independent} 章中，我们展示了如何学习那些假设支持集为单个或多个低维子空间或低秩高斯分布的简单分布类。为了进一步简化，不同的（隐藏的）线性或高斯模式假设为正交或独立。\footnote{或者可以很容易地减少到这些理想情况。}正如我们所展示的，对于这些特殊的分布，可以推导出相当简单且有效的学习算法，具有正确性和效率保证。相关算法中操作的几何和统计解释也非常清晰。

在实践中，线性和独立性都是相当理想化的假设，现实世界高维数据的分布很少满足这些假设。然而，假设分布的内在维度相对于数据嵌入的环境空间的维度来说是相当低的，这是合理的。在本章中，我们将展示如何在高维空间中学习更一般的低维分布类，该空间不一定是（分段）线性的。通常，真实数据的分布包含多个模式，例如在图像的情况下，对应于不同类别的对象。这些模式可能在统计上不独立，甚至可能具有不同的内在维度。通常，我们只能访问分布的有限样本。

为了学习这样的分布，我们需要解决几个基本问题：
\begin{itemize}
    \item 我们如何表示或近似表示一个相当一般的低维分布？
    \item 我们如何衡量这些分布的复杂性，以便可以有效地利用它们的低维性进行学习？
    \item 我们如何使学习过程在计算上可行，甚至可扩展，因为环境维度很高，样本数量很大？
\end{itemize} 我们将看到，压缩或降维的基本思想，对于线性/独立的情况已经证明是非常有效的，仍然可以作为开发学习一般低维分布的有效计算模型和方法的一般原则。

由于其理论和实践上的重要性，我们将更深入地研究通过压缩学习分布的一般框架，当感兴趣的分布可以很好地建模或近似为低维子空间或低秩高斯分布的混合时，该框架如何得以实现。
%\yima{我没有触及上述介绍。当本章的其余部分稳定时，我们可以重写它。但应该明确说明目的：一个计算框架，力求实现可实现和可实现的压缩编码和解码方案！}

\section{压缩和熵最小化}

\subsection{编码率和熵}
在第 \ref{ch:intro} 章中，我们提到学习的目标是找到生成给定数据集的最简单方法。概念上，Kolmogorov 复杂度旨在提供这种复杂度的度量，但它不可计算，也不与任何实际可以再现数据的可实现方案相关联。因此，我们需要一个替代的、可计算的、可实现的复杂度度量。这使我们想到了 Shannon 在 1948 年引入的{\em 熵}概念 \cite{Shannon-1948}。

为了说明熵的构造性质，我们从最简单的情况开始。假设我们有一个离散随机变量，它取 $N$ 个不同的值，或称为 \textit{符号}，$\{\x_1, \ldots, \x_N\}$，每个值的概率相等，即 $1/N$。那么我们可以使用 $i$ 的 $\log_2 N$ 位二进制表示来编码每个符号 \(\vx_{i}\)。这种编码方案可以推广到编码任意离散分布 \cite{Cover-Thomas}：给定一个分布 \(p\)，使得 $\sum_{i=1}^N p(\x_i) = 1$，我们可以将每个符号 $\x_i$ 与概率 $p(\x_i)$ 对应的二进制码的大小设为 $\log_2 [1/p(\x_i)] = - \log_2 p(\x_i)$ 位。因此，编码分布 $p(\cdot)$ 中的任何样本所需的平均位数，即 \textit{编码率}，由以下表达式给出：\footnote{根据信息论的惯例 \cite{Cover-Thomas}，这里的 $\log$ 以 $2$ 为底。因此，熵以（二进制）位为单位。}
\begin{equation}
       H(\x) \doteq \mathbb{E}[\log 1/p(\x)]  = - \sum_{i=1}^N p(\x_i) \log  p(\x_i).
       \label{eqn:entropy-discrete}
\end{equation}
这被称为（离散）分布 $p(\cdot)$ 的 \textit{熵}。注意，这个熵总是非负的，并且仅当 $p(\x_i) = 1$ 对于某个 $\x_i$ 且 $i \in [N]$ 时为零。\footnote{这里我们使用了 $\lim_{p\rightarrow 0} p \log p = 0$ 的事实。}

\subsection{微分熵}

当随机变量 $\x$ 是连续的，并且具有概率密度 $p(\x)$ 时，我们可以将上述和式 \eqref{eqn:entropy-discrete} 的极限视为一个积分：\footnote{更精确地说，给定一个连续变量 $\x$，我们可以用量化大小 $\epsilon > 0$ 对其进行量化。记量化后的离散变量为 $\x^\epsilon$。然后我们可以证明 $H(x^\epsilon) + \log(\epsilon) \approx h(x)$。因此，当 $\epsilon$ 很小时，微分熵 $h(x)$ 可能为负 \cite{Cover-Thomas}。}
\begin{equation}
h(\x) \doteq  - \int p(\x) \log p(\x) ~d\x.
\label{eqn:entropy-differential}
\end{equation}

\begin{example}[高斯分布的熵]
通过直接计算，可以证明高斯分布 $x \sim \mathcal{N}(\mu, \sigma^2)$ 的熵为：
\begin{equation}
    h(x) = \frac{1}{2}\log (2\pi \sigma^2) + \frac{1}{2}.
    \label{eqn:entropy-Gaussian}
\end{equation}
还已知高斯分布是所有具有相同方差 $\sigma$ 的分布中熵最大的。多变量高斯分布 $\x \sim \mathcal{N}(\boldsymbol{\mu}, \boldsymbol{\Sigma})$ 在 $\mathbb{R}^D$ 中的熵为：
\begin{equation}
    h(\x) = \frac{D}{2}(1 + \log(2\pi)) + \frac{1}{2}\log\det(\boldsymbol{\Sigma}).
    \label{eqn:entropy-Gaussian-multi}
\end{equation}
\end{example}

与离散分布的熵类似，我们希望微分熵与某种可实现的编码方案的编码率相关联。例如，我们可以用大小为 $\epsilon >0$ 的网格对分布的域进行离散化。离散化后分布的编码率可以视为微分熵的一个近似 \cite{Cover-Thomas}。

需要注意的是，微分熵的定义存在一些注意事项。对于高维空间中的分布，当其支撑集退化（低维）时，其微分熵可能没有定义。即使对于高斯分布 \eqref{eqn:entropy-Gaussian-multi}，当 \(\vSigma\) 奇异时，$h(\x) = -\infty$。在这种情况下，如何正确量化或编码这样的分布并不明显。然而，退化的（高斯）分布正是高维空间中低维分布的最简单且可能最重要的实例。在本章中，我们将讨论一个完全解决这种退化问题的方法。

%\sdb{看起来这并不那么直接（例如，大小为$N$的字母表上的均匀分布的熵为$log_2 N$比特）...我认为正确的方法是考虑将连续随机变量$\vx$的值空间量化为$N$个桶，计算结果的离散熵，对离散熵进行一定的修正，然后取$N\to\infty$...这个修正项似乎允许微分熵为负，而离散熵总是非负的...}
%\sdb{对此不太确定...例如，高斯分布不是相对于适当归一化（单位方差）的分布来说是“最大熵”的吗？但我的担忧更广泛（根据这个标准，所有完美的低维分布的熵都是$-\infty$，比如）...}

\subsection{最小化编码率}\label{sub:min_entropy}
记住，学习问题是从一组从分布中抽取的样本$\X = \{\x_1, \ldots, \x_N\}$中恢复一个（可能是连续的）分布$p(\x)$。鉴于这里感兴趣的分布是（几乎）低维的，我们应该期望它们的（微分）熵非常小。但与我们在前一章中研究的情况不同，通常我们不知道分布$p(\x)$所属的（解析）低维模型族。因此，检查熵是否很小似乎是唯一可以依赖的指导方针，以识别和建模分布。

现在仅给出样本而不了解 $p(\x)$ 是什么，理论上它们可以被解释为来自任何通用分布的样本。特别是，它们可以被解释为以下任何一种情况：
\begin{enumerate}
    \item 作为经验分布 $\X$ 本身的样本，每个样本 $\x_i, i=1, \ldots, N$ 的概率为 $1/N$。
    \item 作为标准正态分布 $\x^n \sim \mathcal{N}(\boldsymbol{0}, \sigma^2 \boldsymbol{I})$ 的样本，方差 $\sigma^2$ 足够大（例如大于样本范数）；
    \item 作为正态分布 $\x^e \sim \mathcal{N}(\boldsymbol{0}, \hat{\boldsymbol{\Sigma}})$ 的样本，协方差 $\hat{ \boldsymbol{\Sigma}} = \frac{1}{N} \X \X^T$ 是样本的经验协方差；
    \item 作为分布 $\hat \x \sim q(\hat \x)$ 的样本，该分布近似于真实分布 $p(\x)$。
\end{enumerate}
现在的问题是，哪一种更好，从什么意义上说？假设你认为这些数据 $\X$ 是从某个特定分布 $q(\x)$ 中抽取的，该分布可能是上述考虑的分布之一。然后我们可以使用分布 $q(\x)$ 的最优码本对数据点进行编码。所需的平均编码长度（或编码率）由下式给出：
\begin{equation}
    \frac{1}{N}\sum_{i=1}^N -\log q(\x_i) \quad \approx \quad - \int p(\x) \log q(\x) ~d\x 
\end{equation}
当样本数量 $N$ 变大时。如果我们已经确定了正确的分布 $p(\x)$，编码率由熵 $- \int p(\x) \log p(\x) ~d\x$ 给出。事实证明，上述编码长度 $- \int p(\x) \log q(\x) ~d\x$ 除非 $q(\x) = p(\x)$，否则总是大于或等于熵。它们之间的差异，记为
\begin{eqnarray}
    D(p \| q) &\doteq& - \int p(\x) \log q(\x) ~d\x - \Big(- \int p(\x) \log p(\x) ~d\x \Big)\\
    &=& \int p(\x) \log \frac{p(\x)}{q(\x)} d\x 
\end{eqnarray}
称为 {\em Kullback-Leibler} 散度，或相对熵。可以证明这个量总是非负的。
\begin{theorem}[信息不等式]\label{thm:information-inequality}
设 $p(\x), q(\x)$ 是两个概率密度函数（具有相同的支撑）。则 $D(p \| q) \ge 0$，当且仅当 $p(\x) = q(\x)$ 时等号成立。
\end{theorem}
\begin{proof}
    \begin{eqnarray*}
    - D(p \| q) 
    &=& - \int p(\x) \log \frac{p(\x)}{q(\x)} d\x 
    =  \int p(\x) \log \frac{q(\x)}{p(\x)} d\x \\
    &\le& \log \int p(\x)  \frac{q(\x)}{p(\x)} d\x \label{eqn:jensen-KL} 
    = \log \int q(\x) d\x = \log 1 = 0,
    \end{eqnarray*}
其中第一个不等式来自 {\em Jensen's inequality} 和函数 $\log(\cdot)$ 严格凹的事实。等号成立当且仅当 $p(\x) = q(\x)$。
\end{proof}

因此，给定一组采样数据 $\X$，为了确定 $p(\x^n)$、$p(\x^e)$ 和 $p(\hat \x)$ 中哪个更好，我们可以比较它们对 $\X$ 的编码率，看看哪个编码率最低。我们从上面知道，一个分布的（理论上可实现的）编码率与其熵密切相关。一般来说，我们有：
\begin{equation}
    h(\x^n) > h(\x^e) > h(\hat \x).
\end{equation}
因此，如果数据 $\X$ 通过与这些分布相关的码本进行编码，编码率在一般情况下会按以下顺序递减：
\begin{equation}
    p(\x^n) \rightarrow p(\x^e) \rightarrow p(\hat \x).
\end{equation}

这一观察为我们提供了一个关于如何追求具有低维结构的分布 \(p(\x)\) 的一般指导方针。它提出了两种可能的方法：
\begin{enumerate}
    \item 从一个具有高熵的一般分布（例如正态分布）开始，通过减少熵逐步将分布转换为数据的（经验）分布。
    \item 在一个大类的（参数或非参数）分布中，这些分布具有明确的编码方案，可以编码给定的数据，逐步寻找更好的编码方案，以获得更低的编码率。
\end{enumerate}
%\sdb{为什么不通过信息增益的角度来呈现这些呢？我认为我理解了这里的论点（一种“从粗到细”的分布模型进展，模型按熵分级），但它并没有告诉我们“何时停止”沿着这个分级层次继续下去。直觉上，似乎应该在模型很好地解释我们的数据时“停止”，那么问题是如何衡量这一点以及如何使其计算化...}
概念上，这两种方法本质上是在做同样的事情。对于第一种方法，我们需要确保这样的转换路径存在并且可以计算。对于第二种方法，需要确保所选的分布族足够丰富，能够紧密逼近（或包含）真实分布。对于这两种方法，我们需要确保具有较低熵或更好编码率的解可以高效计算，并且快速收敛到所需的分布。\footnote{例如，现实世界数据（如图像和文本）的分布。} 我们将在本章的剩余部分中探讨这两种方法。  %\sdb{这里存在一些紧张关系/未说出口的内容（从 Yi 的幻灯片“房间里的大象”）：我们实际上无法“知道”或甚至在高维空间中从有限样本中高效估计这些熵/散度（例如，前一页的 KL 散度有限和 \(\approx\) ——这不具有维数灾难吗？... 或许这可以成为下一节过渡的“转折点”？（去噪提供了一种在这里进行近似计算的可行方法），或者重新组织...}

\section{通过去噪进行压缩}\label{sub:compression_denoising}

%\yima{这一节需要一些结构... 需要一些自然的介绍和对方法的概念性论证。也许从 Sam 的笔记中关于去噪时熵减少的事实开始？尽管技术事实可以放在附录中供参考。}

%\DP{@YM：您可能希望将以下部分或全部文本移到第 1 章或第 2 章。我无法决定放在哪里最好；您已经在第 1 章中从稍微不同的角度详细介绍了这一点。绝对可能将其放在第 2 章的开头，但它也与您提出的大纲正交。}

在本节中，我们将描述一种通过学习分布的参数编码来学习分布 \(p(\vx)\) 的 \textit{自然} 和 \textit{计算上可行} 的方法，使得表示具有最小的熵或编码率。在这一过程中，我们将展示如何使用这种编码将标准高斯分布的高熵样本转换为目标分布的低熵样本，如图 \ref{fig:diffusion-chapter3} 所示。这提供了一种利用上述两种方法来学习和采样分布的方法。

\begin{figure}[t]
    \centering
\includegraphics[width=\linewidth]{figures/diffusion_pipeline-2.pdf}
    \caption{从高斯分布开始，通过迭代去噪过程收敛到任意数据分布的示意图。}
    \label{fig:diffusion-chapter3}
\end{figure}

\subsection{通过去噪最小化熵}

我们首先希望找到一种方法，将给定的非常嘈杂的样本的熵减少到数据分布中的低熵样本。在这里，我们描述了一种可能的方法——众多方法之一，但可能是解决这个问题最自然的方式。首先，我们找到一种方法来\textit{增加}现有样本的熵。然后，我们找到这个操作的近似逆操作。但一般来说，增加熵的操作没有逆操作，因为原始分布的信息可能会被破坏。因此，我们将处理一个特殊情况，即（1）增加熵的操作具有简单、可计算且可逆的形式；（2）我们可以获得数据分布的（参数化）编码，如上述两个假设中所述。正如我们将看到的，上述两个因素将确保我们所描述的方法是可行的。

我们将以最简单的方式增加熵，即\textit{添加各向同性高斯噪声}。更准确地说，假设\(\vz_{o} \in \R^{d}\)是我们想要估计的随机变量（即\textit{植物}或\textit{真实值}）。假设\(\vz_{\sigma}\)是\(\vz_{o}\)通过加性高斯噪声的扰动，即     \begin{equation}\label{eq:additive_gaussian_noise_model}
    \vz_{\sigma} = \vz_{o} + \sigma \vg, \qquad \vg \sim \dNorm(\vzero, \vI).
\end{equation}
以下定理表明，通过添加高斯噪声，结果分布的微分熵确实增加了。 \sdb{TODO: 使过渡更加流畅。
解决负无穷微分悖论的方法是稍后的有损编码/率失真方案。考虑一下过渡。同时参考现代实践中的VQ-VAE。对于这一部分，如果这是“在非常干净的设置下如何做到的”（稍后通过离散化实现实际应用），只要过渡/上下文良好，这是可以接受的}
\begin{theorem}[扩散与熵]\label{thm:diffusion-and-entropy}
    假设\(\vz_{o} \in \R^{d}\)是一个具有明确定义的微分熵\(h(\vz_o)\)的随机变量，%
    %with \(\Ex\norm{\vz_{o}}_{2}^{2} < \infty\),
    \(\vg \sim \dNorm(\vzero, \vI)\)是一个标准高斯随机变量，独立于\(\vz_{o}\)，\(\sigma > 0\)是一个正实数，且\(\vz_{\sigma} = \vz_{o} + \sigma \vg\)。那么\(h(\vz_{\sigma}) > h(\vz_{o})\)。
\end{theorem}
证明简洁但需要一些较为复杂的数学知识，超出了本书的范围。我们在附录\ref{app:diffusion-denoising}中提供了详细信息。
证明通过极限论证，严格地证明了\(h(\vz_o) = -\infty\)对于大量退化分布成立，正如我们之前所声称的，因此在\cref{thm:diffusion-and-entropy}中所断言的熵增加也适用于这一重要类别的分布。
% \yima{可能需要陈述为定理或命题。}

逆向操作去除各向同性高斯噪声的过程称为\textit{去噪}。这是信号处理和系统理论中的一个经典且研究充分的主题，例如维纳滤波器和卡尔曼滤波器。第\ref{ch:classic}章中讨论的几个问题，如PCA、ICA和字典学习，都是去噪问题的具体实例。具体来说，去噪问题可以表述为尝试学习一个函数\(\vmu_{\sigma} \colon \R^{d} \to \R^{d}\)，该函数解决最大似然估计问题，以从\(\vmu_{\sigma}\)中估计\(\vz_{o}\)：
\begin{equation}\label{eq:denoiser_maximum_likelihood}
    \vmu_{\sigma} \in \argmax_{\substack{f \colon \R^{d} \to \R^{d},\   
    %\\%f\ \text{可测} \\ 
    \Ex\norm{f(\vz_{\sigma})}_{2}^{2} < \infty}}\Ex \log p_{\sigma \mid o}(f(\vz_{\sigma})),
\end{equation}
特别是，去噪器\(\vmu_{\sigma}\)找到使噪声观测最有可能的真实数据值。很容易证明，找到函数\(\vmu_{\sigma}\)，在高斯噪声模型\eqref{eq:additive_gaussian_noise_model}下，等价于解决以下问题
%\footnote{如果你不确定\textit{可测}是什么意思，不要担心；这只是技术条件，对于任何你可以写出的函数都成立。其确切定义可以在任何研究生水平的实分析或概率论教科书中查到。}
\begin{equation}\label{eq:denoiser_conditional_expectation}
    \vmu_{\sigma} \in \argmin_{\substack{f \colon \R^{d} \to \R^{d},\ %\\% \text{\(f\) 可测} \\ 
    \Ex\norm{f(\vz_{\sigma})}_{2}^{2} < \infty}}\Ex\norm{\vz_{o} - f(\vz_{\sigma})}_{2}^{2}.
\end{equation}
也就是说，对于某个包含一些潜在去噪函数的函数类\(\cF\)，例如具有相同架构和不同权重的神经网络族，我们可以解决去噪问题
\begin{equation}\label{eq:denoiser_learning_problem}
    f^{\star} \in \argmin_{f \in \cF}\Ex\norm{\vz_{o} - f(\vz_{\sigma})}_{2}^{2}.
\end{equation}
% 事实证明，解决\eqref{eq:denoiser_conditional_expectation}中的最小化问题比解决最大似然问题\eqref{eq:denoiser_maximum_likelihood}要\textit{容易}得多。
为了解决这个问题，我们使用以下概述的算法\ref{alg:denoiser}。
\begin{algorithm}
    \label{alg:denoiser}
    \caption{学习去噪器}
    \begin{algorithmic}[1]
        \Require{一个优化算法，一个参数化函数类\(\cF = \{f_{\vtheta} \colon \vtheta \in \vTheta\}\)，一个初始参数\(\vtheta^{1}\)，一个噪声水平\(\sigma > 0\)，以及数据\(\vz_{o}^{1}, \dots, \vz_{o}^{N} \simiid \vz_{o}\)。}
        \Ensure{一个去噪器\(f_{\vtheta^{\star}}\)。}

\Procedure{TrainDenoiser}{\texttt{OptUpdate}, $\cF$, $\vtheta^{1}$, $\sigma$, $\{\vz_{o}^{i}\}_{i = 1}^{N}$}
        \For{每个迭代\(j \in [M]\)}
            \For{每个样本\(i \in [N]\)}
                \State{\(\vg^{i, j} \gets \) 从\(\dNorm(\vzero, \vI)\)中采样} \Comment{\(\vg^{i, j}\)与所有其他变量独立}
                \State{\(\vz_{\sigma}^{ij} \gets \displaystyle \vz_{o}^{i} + \sigma \vg^{i, j}\)}
            \EndFor
            \State{\(\ell^{j} \gets \displaystyle\frac{1}{N}\sum_{i = 1}^{N}\norm{\vz_{o}^{i} - f_{\vtheta^{j}}(\vz_{\sigma, j}^{i})}_{2}^{2}\)} \Comment{\(\ell^{j}\)是“真实”去噪损失的估计}
            \State{\(\vtheta^{j + 1} \gets \displaystyle \texttt{OptUpdate}(\vtheta^{j}; \ell^{j})\)} \Comment{通过优化更新\(\vtheta^{j + 1}\)，即\(\vtheta^{j + 1} \gets \vtheta^{j} - \alpha \nabla_{\vtheta}\ell^{j}\)用于SGD。}
            \EndFor
        \State{\Return{\(f_{\vtheta^{M+1}}\)}}
        \EndProcedure
    \end{algorithmic}
\end{algorithm}

算法中唯一需要指定的是函数类\(\mathcal{F}\)的形式。对于特殊的分布类，如高斯分布，函数形式可以显式导出。

\begin{example}[从高斯分布中去噪高斯噪声]\label{example:gaussian_denoising}
    考虑 \(\vz_{o}\) 是一个高斯分布，例如 \(\vz_{o} \sim \dNorm(\vu, \vSigma)\)。那么 \(\vz_{\sigma} \sim \dNorm(\vu, \vSigma + \sigma^{2}\vI)\)，其中 \(\vSigma\) 是非奇异的。那么 \(\vz_{o}\) 和 \(\vz_{\sigma}\) 联合高斯分布，即
    \begin{equation}
        \mat{\vz_{o} \\ \vz_{\sigma}} \sim \dNorm\rp{\mat{\vu \\ \vu}, \mat{\vSigma & \vSigma \\ \vSigma & \vSigma + \sigma^{2}\vI}},
    \end{equation}
    并且给定 \(\vz_{\sigma}\) 的 \(\vz_{o}\) 的后验分布是（著名地）
    \begin{equation}
        \vz_{o} \mid \vz_{\sigma} \sim \dNorm(\vu + \vSigma(\vSigma + \sigma^{2}\vI)^{-1}(\vz_{\sigma} - \vu), \vSigma - \vSigma (\vSigma + \sigma^{2}\vI)^{-1}\vSigma).
    \end{equation}
    因此我们有
    \begin{equation}\label{eq:gaussian_denoiser}
        \vmu_{\sigma}(\vxi) = \vu + \vSigma(\vSigma + \sigma^{2}\vI)^{-1}(\vxi - \vu),
    \end{equation}
    这为两个去噪问题提供了答案。令 \(\vSigma = \vV\vLambda\vV^{\top}\)，其中 \(\vLambda = \diag(\lambda_{1}, \dots, \lambda_{d})\)，我们有
    \begin{equation}
        \vmu_{\sigma}(\vxi) = \vu + \vV\diag\mat{\lambda_{1}/(\lambda_{1} + \sigma^{2}) & & \\ & \ddots & \\ & & \lambda_{d}/(\lambda_{d} + \sigma^{2})}\vV^{\top}(\vxi - \vu).
    \end{equation}
    这表明随着 \(\sigma\) 的增长，去噪器在 \(\vSigma\) 的特征值最小的方向上最快地将残差 \(\vxi - \vu\) 压缩到 \(0\)。此外，如果我们试图学习一个高斯分布（例如），我们可以用均值 \(\vu\) 和协方差 \(\vSigma\) 来参数化去噪器，使得 \(\vtheta = (\vu, \vSigma)\)，并考虑形式为 \eqref{eq:gaussian_denoiser} 的去噪器。
\end{example}

根据 \eqref{eq:denoiser_conditional_expectation} 的定义，可以证明 \(\vmu_{\sigma}\) 是给定 \(\vz_{\sigma}\) 的 \(\vz_{o}\) 的条件期望，\footnote{我们将其留给读者验证。} 即
\begin{equation}
    \vmu_{\sigma}(\vxi) = \Ex[\vz_{o} \mid \vz_{\sigma} = \vxi].
\end{equation}
可以利用这一点来证明去噪减少了熵 \(h(\vz)\)，如所期望的那样。
\begin{theorem}[去噪和熵]\label{thm:conditioning_reduces_entropy}
    假设 \(\vz_{o} \in \R^{d}\) 是一个随机变量，
    %with \(\Ex\norm{\vz_{o}}_{2}^{2} < \infty\),
    \(\vg \sim \dNorm(\vzero, \vI)\) 是一个标准高斯随机变量
    与 \(\vz_{o}\) 独立，\(\sigma > 0\) 是一个正实数，
    \(\vz_{\sigma} = \vz_{o} + \sigma \vg\)，并且 \(\vmu_{\sigma}(\vz)
    = \Ex[\vz_{o} \mid \vz_{\sigma} = \vz]\)。在某些关于 \(\vz_o\) 的正则性条件下（详见 \Cref{app:diffusion-denoising}），有
    \(h(\vmu_{\sigma}(\vz_{\sigma})) < h(\vz_{\sigma})\)。
\end{theorem}
证明稍显技术性，详见 \Cref{app:diffusion-denoising}。关于 \(\vz_o\) 的必要正则性条件，假设 \(\vz_o\) 是任何随机变量，其中添加了任意小量的独立高斯噪声就足够了。在后续内容中，我们将使用 \textit{迭代去噪} 来降低高熵分布的熵，直到它达到数据分布，实现 \Cref{sub:min_entropy} 结尾处讨论的追求分布的第二种方法：这使得 \Cref{thm:conditioning_reduces_entropy} 适用于迭代去噪过程中的各个阶段，对应于从数据分布中去除先前添加的噪声。

到目前为止，我们已经讨论了使用特定定义的“学习”来学习分布。然而，在文献中，通常情况下，“学习分布”对应于学习数据分布的概率测度或密度。我们现在讨论我们的去噪公式与这种以密度为中心的范式中学习分布的直接联系。这种联系源于 \textit{Tweedie's formula} \cite{Robbins1956AnEB}，我们重复如下。

\begin{theorem}[Tweedie's Formula]\label{thm:tweedie}
    假设 \(\vz_{o} \in \R^{d}\) 是一个随机变量，满足 \(\Ex\norm{\vz_{o}}_{2}^{2} < \infty\)，\(\vg \sim \dNorm(\vzero, \vI)\) 是一个标准高斯随机变量，且与 \(\vz_{o}\) 独立，\(\sigma > 0\) 是一个正实数，\(\vz_{\sigma} = \vz_{o} + \sigma \vg\)，并且 \(p_{\sigma}\) 是 \(\vz_{\sigma}\) 的概率密度函数。那么 
    \begin{equation}
        \Ex[\vz_{o} \mid \vz_{\sigma} = \vxi] = \vxi + \sigma^{2}\nabla \log p_{\sigma}(\vxi).
    \end{equation}
\end{theorem}
\begin{proof}
    为了证明这个定理，假设 \(\vz_{o}\) 的概率密度函数为 \(p_{o}\)（尽管即使 \(\vz_{o}\) 没有密度函数，结果仍然成立）。那么，注意到 \(p_{\sigma}\) 可以分解为 
    \begin{equation}
        p_{\sigma}(\vxi) = \int_{\R^{d}}p_{o}(\vx)p_{\sigma \mid o}(\vxi \mid \vx)\odif{\vx},
    \end{equation}
    其中 \(p_{\sigma \mid o}\) 是给定 \(\vz_{o}\) 的观测值时 \(\vz_{\sigma}\) 的概率密度。相应地定义 \(p_{o \mid \sigma}\)。现在，根据噪声模型，我们有 
    \begin{equation}
        p_{\sigma \mid o}(\vxi \mid \vx) = \phi(\vxi \mid \vx, \sigma^{2}\vI),
    \end{equation}
    其中 \(\phi(\vmu, \vSigma)\) 是均值为 \(\vmu\) 和协方差为 \(\vSigma\) 的正态分布的密度函数。那么
    \begin{align}
        \nabla \log p_{\sigma}(\vxi)
        &= \frac{1}{p_{\sigma}(\vxi)}\nabla p_{\sigma}(\vxi) \\ 
        &= \frac{1}{p_{\sigma}(\vxi)}\nabla_{\vxi}\int_{\R^{d}}p_{o}(\vx)\phi(\vxi \mid \vx, \sigma^{2}\vI)\odif{\vx} \\ 
        &= \frac{1}{p_{\sigma}(\vxi)}\int_{\R^{d}}p_{o}(\vx)\bs{\nabla_{\vxi}\phi(\vxi \mid \vx, \sigma^{2}\vI)}\odif{\vx} \\ 
        &= \frac{1}{p_{\sigma}(\vxi)}\int_{\R^{d}}p_{o}(\vx)\bs{\phi(\vxi \mid \vx, \sigma^{2}\vI)\frac{\vx - \vxi}{\sigma^{2}}}\odif{\vx} \\
        &= \frac{1}{\sigma^{2}p_{\sigma}(\vxi)}\int_{\R^{d}}p_{o}(\vx)\phi(\vxi \mid \vx, \sigma^{2}\vI)[\vx - \vxi]\odif{\vx} \\
        &= \frac{1}{\sigma^{2}p_{\sigma}(\vxi)}\int_{\R^{d}}p_{o}(\vx)\phi(\vxi \mid \vx, \sigma^{2}\vI)\vx\odif{\vx} \\ 
        &\qquad - \frac{1}{\sigma^{2}p_{\sigma}(\vxi)}\vxi\int_{\R^{d}}p_{o}(\vx)\phi(\vxi \mid \vx, \sigma^{2}\vI)\odif{\vx} \\
        &= \frac{1}{\sigma^{2}p_{\sigma}(\vxi)}\int_{\R^{d}}p_{o}(\vx)\phi(\vxi \mid \vx, \sigma^{2}\vI)\vx\odif{\vx} - \frac{1}{\sigma^{2}p_{\sigma}(\vxi)}\vxi p_{\sigma}(\vxi), \\
        &= \frac{1}{\sigma^{2}p_{\sigma}(\vxi)}\int_{\R^{d}}p_{o}(\vx)\phi(\vxi \mid \vx, \sigma^{2}\vI)\vx\odif{\vx} - \frac{1}{\sigma^{2}}\vxi.
    \end{align}
    我们来考虑剩余的积分。记住 \(\phi(\vxi \mid \vx, \sigma^{2}\vI) = p_{\sigma \mid o}(\vxi \mid \vx)\)。根据贝叶斯定理，我们有 
    \begin{equation}
        p_{o}(\vx)\phi(\vxi \mid \vx, \sigma^{2}\vI) = p_{o}(\vx)p_{\sigma \mid o}(\vxi \mid \vx) = p_{\sigma}(\vxi)p_{o \mid \sigma}(\vx \mid \vxi), 
    \end{equation}
    这使得我们可以将被积函数改写为更有利于计算的形式，最终得到 
    \begin{align}
        \nabla \log p_{\sigma}(\vxi)
        &= \frac{1}{\sigma^{2}p_{\sigma}(\vxi)}\int_{\R^{d}}p_{o}(\vx)\phi(\vxi \mid \vx, \sigma^{2}\vI)\vx\odif{\vx} - \frac{1}{\sigma^{2}}\vxi \\ 
        &= \frac{1}{\sigma^{2}p_{\sigma}(\vxi)}\int_{\R^{d}}p_{\sigma}(\vxi)p_{o \mid \sigma}(\vx \mid \vxi)\vx\odif{\vx} - \frac{1}{\sigma^{2}}\vxi \\ 
        &= \frac{1}{\sigma^{2}}\int_{\R^{d}}p_{o \mid \sigma}(\vx \mid \vxi)\vx\odif{\vx} - \frac{1}{\sigma^{2}}\vxi \\ 
        &= \frac{1}{\sigma^{2}}\bp{\Ex[\vz_{o} \mid \vz_{\sigma} = \vxi] - \vxi}.
    \end{align}
    重新排列等式即可证明定理。
\end{proof}
%\yima{Tweedie's formula for the scalar case... as an example of an exercise.}
%\DP{Exercise \ref{exercise:vp_tweedie_denoising} gives a better exercise about these contents. Scalar Tweedie's formula is not a good question since it just repeats the same calculation as here.}

这个结果建立了估计条件期望 \(\Ex[\vz_{o} \mid \vz_{\sigma}]\) 和估计数据分布 \(\nabla \log p_{\sigma}\)（即梯度对数密度）之间的联系，后者也被称为 \textit{(Hyv{\"a}rinen) 分数 \cite{hyvarinen05a}}。我们已经有一个估计 \(\Ex[\vz_{0} \mid \vz_{\sigma}]\) 的方法，因此我们也有一个估计 \(\nabla \log p_{\sigma}\) 的方法。正如我们在第 \ref{ch:intro} 章中所看到的，图 \ref{fig:score-function} 说明了在 1D 分布情况下分数的几何直觉。接下来，我们将讨论如何使用我们对 \(\nabla \log p_{\sigma}\) 的估计来逐步将高熵分布转换为（低维）数据分布。

% \yima{如何解决这个问题？是否需要参数化 \(f\)？读者需要为每个 \(\z_o\) 解决这个问题吗？}

%\yima{这一段。请不要像读者已经知道去噪问题的答案一样发表评论。这是一本书，不是论文。我们应该直接描述如何解决去噪问题。在事实呈现之后可以进行一些评论。}
%在非常高的层面上，去噪遵循我们声明的学习分布的一般方法。然而，条件期望，即最优去噪器，与数据分布（及其扰动版本）之间也存在非常具体的联系。这些联系大多源于 \textit{Tweedie's formula} \DP{引用: Robbins 1957，或更早？} \DP{引用: Jong-Chul Ye DPM，这是我所知道的最一般的Tweedie's formula形式} \DP{引用: Druv的MAE论文}，其最简单的形式如下。考虑模型 \eqref{eq:additive_gaussian_noise_model}，设 \(p_{\sigma}\) 是 \(\vz_\sigma\) 的概率密度，设 \(\vs_{\sigma} \doteq \nabla \log p_{\sigma}\)。那么
% \begin{equation}\label{eq:tweedie_simple}
%     \vmu_{\sigma}(\vz) \doteq \Ex[\vz_{o} \mid \vz_{\sigma} = \vz] = \vz + \sigma^{2}\nabla \log p_{\sigma}(\vz) \doteq \vz + \sigma^{2}\vs_{\sigma}(\vz).
% \end{equation}
% 因此，学习数据分布的扰动版本通过 \(\nabla \log p^{\sigma}\)，也称为（Hyvarinen）分数 \DP{引用: aapo score matching} \DP{引用: sohl-dickstein, ganguli}，与解决简单的条件期望回归问题 \eqref{eq:denoiser_learning_problem} 是等价的，后者更容易处理。 \yima{这句话很别扭：什么是等价的？哪个更易处理？我们需要确切地描述条件期望如何解决，分数函数如何表示和估计，至少在高斯噪声的情况下。清楚地描述一种方法，并可能评论其他替代方法来估计分数等。或者留作练习。}

% 可以证明，通过上述去噪（使用分数）得到的数据分布的熵总是减少的：$$h(\boldsymbol{\mu}_\sigma) < h(\z_\sigma).$$
% 再次，参见附录 \ref{app:diffusion-denoising} 以获取证明。 \yima{可能作为定理陈述。}

\subsection{通过迭代去噪进行采样}

%\yima{也许添加一个关于如何将这应用于单个子空间/高斯的例子，即PCA的情况？也许在本小节末尾，因为下一小节处理更一般的高斯混合。那么本节将与第3.3节的流程平行。}

之前，我们已经讨论了一个通过拟合参数函数族中的去噪器来学习数据分布的参数编码的方法。我们还在定理 \ref{thm:conditioning_reduces_entropy} 中观察到，应用去噪器会减少噪声分布的熵。在本节中，我们将讨论如何仔细使用这个去噪器来逐步减少噪声分布的熵，直到达到数据分布，从而 \textit{采样} 该分布。这个过程称为 \textit{迭代去噪} 或 \textit{迭代压缩}，是第 \ref{sub:min_entropy} 节末尾提出的第一个学习或追求分布的方法的具体实现。

首先，注意到我们的去噪器是从 \(\vz_{\sigma}\) 映射到 \(\vz_{o}\) 的估计值。确实有 \(h(\Ex[\vz_{o} \mid \vz_{\sigma}]) < h(\vz_{\sigma})\)。但是，量 \(\Ex[\vz_{o} \mid \vz_{\sigma}]\) 通常不是一个数据分布支持上的点，而是一个这样的点的平均值，因此不是一个真正的数据点。因此，为了仔细地减少熵以便采样，我们 \textit{不能} 只使用一次去噪器。

相反，我们选择将问题分解为 \textit{许多小步骤} —— 一次减少一点噪声（或减少一点熵）（一个计算上更容易的问题），并将结果串联起来，以从大噪声/高熵分布中去噪/减少熵（一个困难的问题）。更精确地说，我们有
\begin{equation}
    \vz_{\sigma} = \vz_{o} + \sigma \vg \implies \vg = \frac{\vz_{\sigma} - \vz_{o}}{\sigma}.
\end{equation}
然后，对于任何 \(\delta > 0\)，我们有
\begin{align}
    \Ex[\vz_{\sigma - \delta} \mid \vz_{\sigma}]
    &= \Ex[\vz_{o} + (\sigma - \delta)\vg \mid \vz_{\sigma}] \\ 
    &= \Ex[\vz_{\sigma} - \delta \vg \mid \vz_{\sigma}] \\ 
    &= \vz_{\sigma} - \delta \Ex[\vg \mid \vz_{\sigma}] \\
    &= \vz_{\sigma} - \delta \Ex\rs{\frac{\vz_{\sigma} - \vz_{o}}{\sigma} \given\ \vz_{\sigma}} \\
    &= \bp{1 - \frac{\delta}{\sigma}}\vz_{\sigma} + \frac{\delta}{\sigma}\Ex[\vz_{o}\ \given\ \vz_{\sigma}].
\end{align}
因此，我们可以进行迭代去噪 —— 重复地从 \(\vz_{\sigma}\) 恢复 \(\vz_{\sigma - \delta}\)，最终我们尝试从 \(\vz_{\sigma}\) 恢复 \(\vz_{o}\)，其中 \(\sigma\) 很小。此时，条件期望向量 \(\Ex[\vz_{o} \mid \vz_{\sigma}]\) 将非常接近数据分布的支持，允许我们近似采样。形式上，我们可以设置以下方法：设置一个终端迭代 \(M\) 和 \textit{噪声计划} \((\sigma^{(m)})_{m = 0}^{M - 1}\)，使得 \(\sigma^{(0)} > \sigma^{(1)} > \cdots > \sigma^{(M - 1)} > 0\)，并且计划中后续噪声水平之间的差异可以写成
\begin{equation}
    \delta^{(m)} \doteq \begin{cases}\sigma^{(m)} - \sigma^{(m + 1)}, & \text{如果}\ 0 \leq m < M - 1 \\ \sigma^{(M - 1)}, & \text{如果}\ m = M - 1.\end{cases}
\end{equation}
然后，我们有迭代
\begin{align}
    \vz^{(0)} 
    &\sim \dNorm(\vzero, (\sigma^{(0)})^{2}\vI) \label{eq:denoising_ve_init}, \\ 
    \vz^{(m + 1)} 
    &= \bp{1 - \frac{\delta^{(m)}}{\sigma^{(m)}}}\vz^{(m)} + \frac{\delta^{(m)}}{\sigma^{(m)}}\Ex[\vz_{o} \mid \vz_{\sigma^{(m)}} = \vz^{(m)}], \label{eq:denoising_ve_iter} \\
    &\qquad  \qquad \forall m \in \{0, 1, \dots, M - 1\}
\end{align}
这种方法在 \(M\) 步中每次去噪 \(\delta^{(m)}\)，使得 \(\vz^{(m)}\) 具有（近似）噪声水平 \(\sigma^{(m)}\)。出于历史原因，得分重述更受欢迎，我们可以使用 Tweedie 公式获得它，即
\begin{align}
    \vz^{(m + 1)}
    &= \bp{1 - \frac{\delta^{(m)}}{\sigma^{(m)}}}\vz^{(m)} + \frac{\delta^{(m)}}{\sigma^{(m)}}\Ex[\vz_{o} \mid \vz_{\sigma^{(m)}} = \vz^{(m)}] \\ 
    &=  \bp{1 - \frac{\delta^{(m)}}{\sigma^{(m)}}}\vz^{(m)} + \frac{\delta^{(m)}}{\sigma^{(m)}}\bs{\vz^{(m)} + (\sigma^{(m)})^{2}\nabla \log p_{\sigma^{(m)}}(\vz^{(m)})} \\ 
    &= \vz^{(m)} + \delta^{(m)} \sigma^{(m)}\nabla \log p_{\sigma^{(m)}}(\vz^{(m)}).
\end{align}
因此，我们将每个去噪步骤重述为在对数密度 \(\log p_{\sigma^{(m)}}(\vz^{(m)})\) 上以步长 \(\delta^{(m)} \sigma^{(m)}\) 进行梯度步。图 \ref{fig:diffusion-chapter3} 以视觉方式描述了这样的梯度步，以及迭代去噪的几何结构。为了方便起见，下面简要地呈现了一个算法。
\begin{algorithm}\label{alg:iterative_denoising}
    \caption{迭代去噪}
    \begin{algorithmic}[1]
        \Require{噪声计划 \((\sigma^{(m)})_{m = 0}^{M - 1}\) 使得 \(\sigma^{(0)} > \sigma^{(1)} > \cdots > \sigma^{(M - 1)} > 0\)，样本 \(\{\vz_{o}^{i}\}_{i = 1}^{N} \simiid \vz_{o}\)。}
        \Ensure{从 \(\vz_{o}\) 的分布中（近似）抽取的样本 \(\vz^{(M)}\)。}

\Procedure{迭代去噪}{$(\sigma^{(m)})_{m = 0}^{M - 1}$, $\{\vz_{o}^{i}\}_{i = 1}^{N}$}
        
        \For{\(m \in \{0, 1, \dots, M-2\}\)} \Comment{计算噪声差异 \(\delta^{(m)}\)}
            \State{\(\delta^{(m)} \gets \sigma^{(m)} - \sigma^{(m + 1)}\)}
        \EndFor
        \State{\(\delta^{(M - 1)} \gets \sigma^{(M-1)}\)}
        \Statex{}

\State{初始化 \(\vz^{(0)} \gets \dNorm(\vzero, \vI)\)} \Comment{初始化去噪}
        \For{\(m \in \{0, 1, \dots, M-1\}\)}
            \State{\(\hat{\vmu}_{\sigma^{(m)}} \gets \displaystyle \textsc{TrainDenoiser}(\sigma^{(m)}, \{\vz_{o}^{i}\}_{i = 1}^{N})\)} \Comment{来自 \Cref{alg:denoiser}}
            \State{\(\hat{\vs}_{\sigma^{(m)}} \gets \displaystyle \frac{\hat{\vmu}_{\sigma^{(m)}}(\vz^{(m)}) - \vz^{(m)}}{(\sigma^{(m)})^{2}}\)} \Comment{使用 Tweedie 公式（\Cref{thm:tweedie}）转换为近似分数}
            \State{\(\vz^{(m + 1)} \gets \displaystyle \vz^{(m)} + \delta^{(m)}\sigma^{(m)}\hat{\vs}_{\sigma^{(m)}}(\vz^{(m)})\)} \Comment{去噪迭代}
        \EndFor

\State{\Return{\(\vz^{(M)}\)}}
        \EndProcedure
    \end{algorithmic}
\end{algorithm}

一些常见的噪声调度选择包括：
\begin{itemize}
    \item \(\sigma^{(m)} = C(M - m)\) 其中 \(C > 0\)，这是“线性噪声调度”。
    \item \(\sigma^{(m)} = C(\e^{\alpha (M - m)} - 1)\) 其中 \(C > 0\) 和 \(\alpha > 0\)，这是“指数噪声调度”。
\end{itemize}
先前的工作\footnote{\cite{karras2022elucidating} 实验性地评估了不同的噪声调度，而 \cite{bao2022analytic} 构造了理论上最优的噪声调度，尽管需要通过蒙特卡洛方法进行估计。} 研究了“最优”噪声调度的概念。

算法 \ref{alg:iterative_denoising} 即使在分数估计的情况下也能收敛到良好的精度（只要估计足够好）。
\begin{theorem}[《\cite{lee2022convergence}》的定理 2.2，简写]\label{thm:distribution-convergence}
    在某些技术条件下，对于 \(\vz_{o}\) 的噪声水平 \((\sigma^{(m)})_{m = 0}^{M - 1}\) 和从样本 \(\{\vz_{o}^{i}\}_{i = 1}^{N}\) 学习到的分数估计 \(\hat{\vs}_{\sigma^{(m)}}\) 的质量，如果 \(\vz^{(M)}\) 是算法 \ref{alg:iterative_denoising} 中采样过程的输出，则 \(\vz^{(M)} \to \vz_{o}\)（在所谓的 \textit{总变差距离}中）当 \(M, N \to \infty\) 时。
\end{theorem}

\begin{remark}
    在数据 \(\vz_{o}\) 具有低维结构（如第 \ref{ch:classic} 章所述）的情况下，\(\vz_{o}\) 的密度可能不存在。无论如何，这里的技术同样适用，因为 \(\vz_{\sigma}\) 对于任何 \(\sigma > 0\) 都具有密度。\footnote{这是一般情况下难以证明的事实，需要满足所谓的 \textit{H\"ormander 条件}，并且证明 H\"ormander 条件意味着密度的存在需要 Malliavin 微积分。请参考任何高级随机微积分教科书以获取详细信息。} 然而，对于小 \(\sigma\)，分数 \(\hat{\vs}_{\sigma}\) 可能变得越来越数值不稳定（在雅可比矩阵病态或甚至非 Lipschitz 连续的意义上）。因此，建议在从低维分布采样时增加 \(M\)。
\end{remark}

\begin{figure}[h]
    \centering
    \includegraphics[width=0.32\textwidth]{figs_chap3/gaussian_mixture_denoising/data.png}
    \hfill
    \includegraphics[width=0.32\textwidth]{figs_chap3/gaussian_mixture_denoising/m_0.png}
    \hfill
    \includegraphics[width=0.32\textwidth]{figs_chap3/gaussian_mixture_denoising/m_460.png}
    
    \includegraphics[width=0.32\textwidth]{figs_chap3/gaussian_mixture_denoising/m_480.png}
    \hfill
    \includegraphics[width=0.32\textwidth]{figs_chap3/gaussian_mixture_denoising/m_490.png}
    \hfill
    \includegraphics[width=0.32\textwidth]{figs_chap3/gaussian_mixture_denoising/m_500.png}
    \caption{\textbf{去噪以从低秩高斯混合中采样。} \textit{左上角：}我们从三个高斯分布开始，颜色分别为灰色、橙色和红色，分别在\(\R^{3}\)中的一个\(2\)-d子空间和两个\(1\)-d子空间上。 \textit{中间上：}我们通过设置\(\sigma^{(0)} = 10\)开始去噪过程\eqref{eq:denoising_ve_iter}，使得\(\vz^{(0)} \sim \dNorm(\vzero, 10\vI)\)，如\eqref{eq:denoising_ve_init}所示。 \textit{右上角：}经过大部分迭代后，当前迭代\(\vz^{(m)}\)是真实分布的噪声版本。 \textit{左下角，中间下：}带有几个去噪迭代的相同图形。 \textit{右下角：}采样过程的输出，与原始数据集非常相似。} \pw{作为练习？}
    \label{fig:denoising-twolines-oneplane}
\end{figure}

\begin{remark}
    这个“迭代去噪”过程\textit{完全}是扩散模型中使用的采样算法。 （专家或好奇者可能知道这是“方差爆炸概率流”过程。）算法\eqref{alg:iterative_denoising}中的采样过程虽然是确定性的，但等价于一系列可以用于从数据分布中采样的确定性和随机过程。 \footnote{一些例子：\cite{montanari2023sampling}和\cite{song2020denoising}。   %\DP{只是保留这个，不需要在最终版本中出现。}
    } 图\ref{fig:denoising-twolines-oneplane}显示了一个通过迭代去噪高斯分布的随机样本来采样三个退化高斯分布（两条线和一个平面）的示例。
\end{remark}

\begin{example}[从高斯分布中采样]
    在示例\ref{example:gaussian_denoising}中，我们详细讨论了高斯地面真实分布的最优去噪器。上面，我们展示了去噪和采样之间的紧密联系。在这里，我们将计算高斯分布的得分函数和去噪迭代，并与PCA（参见第\ref{ch:classic}章）建立联系。
\end{example}

更具体地说，设 \(\vz_{o} \sim \dNorm(\vu, \vSigma)\)，并设 \(\vz_{\sigma} = \vz_{o} + \sigma \vg\)。我们在示例 \ref{example:gaussian_denoising} 中展示了
    \begin{equation}
        \Ex[\vz_{o} \mid \vz_{\sigma} = \vxi] = \vmu_{\sigma}(\vxi) = \vu + \vSigma(\vSigma + \sigma^{2}\vI)^{-1}(\vxi - \vu).
    \end{equation}
    通过 Tweedie 公式（定理 \ref{thm:tweedie}），我们得到得分：
    \begin{align}
        \nabla \log p_{\sigma}(\vxi)
        &= \frac{\vmu_{\sigma}(\vxi) - \vxi}{\sigma^{2}} \\ 
        &= \frac{1}{\sigma^{2}}\bs{\bp{\vu + \vSigma(\vSigma + \sigma^{2}\vI)^{-1}(\vxi - \vu)} - \vxi} \\
        &= \frac{1}{\sigma^{2}}\bs{\vSigma(\vSigma + \sigma^{2}\vI)^{-1}(\vxi - \vu) - (\vxi - \vu)} \\
        &= \frac{1}{\sigma^{2}}\bs{\vSigma(\vSigma + \sigma^{2}\vI)^{-1} - \vI}(\vxi - \vu).
    \end{align}
    现在，取谱分解 \(\vSigma = \vV\vLambda\vV^{\top}\)，其中 \(\vV \in \O(d)\) 且 \(\vLambda = \diag(\lambda_{1}, \dots, \lambda_{d})\)，其中 \(\lambda_{i} \geq 0\)。我们有
    \begin{align}
        \nabla \log p_{\sigma}(\vxi) &= \frac{1}{\sigma^{2}}\bs{\vV\vLambda\vV^{\top}(\vV\vLambda \vV^{\top} + \sigma^{2}\vI)^{-1} - \vI}(\vxi - \vu) \\
        &= \frac{1}{\sigma^{2}}\bs{\vV\vLambda\vV^{\top}(\vV(\vLambda + \sigma^{2}\vI) \vV^{\top})^{-1} - \vI}(\vxi - \vu) \\
        &= \frac{1}{\sigma^{2}}\bs{\vV\vLambda\vV^{\top}\vV(\vLambda + \sigma^{2}\vI)^{-1} \vV^{\top} - \vI}(\vxi - \vu) \\
        &= \frac{1}{\sigma^{2}}\vV\bs{\vLambda(\vLambda + \sigma^{2}\vI)^{-1} - \vI}\vV^{\top}(\vxi - \vu) \\
        &= \frac{1}{\sigma^{2}}\vV\mat{\lambda_{1}/(\lambda_{1} + \sigma^{2}) - 1 & & \\ & \ddots & \\ & & \lambda_{d}/(\lambda_{d} + \sigma^{2}) - 1}\vV^{\top}(\vxi - \vu) \\
        &= \frac{1}{\sigma^{2}}\vV\mat{-\sigma^{2}/(\lambda_{1} + \sigma^{2}) & & \\ & \ddots & \\ & & -\sigma^{2}/(\lambda_{d} + \sigma^{2})}\vV^{\top}(\vxi - \vu) \\
        &= -\vV\mat{1/(\lambda_{1} + \sigma^{2}) & & \\ & \ddots & \\ & & 1/(\lambda_{d} + \sigma^{2})}\vV^{\top}(\vxi - \vu) \\
        &= -\vV(\vLambda + \sigma^{2}\vI)^{-1}\vV^{\top}(\vxi - \vu) \\
        &= -(\vSigma + \sigma^{2}\vI)^{-1}(\vxi - \vu).
    \end{align}
    
    在这种情况下，得分非常简单，即使去噪器看起来相对复杂。根据算法 \ref{alg:iterative_denoising}，迭代去噪步骤将是
    \begin{equation}
        \vz^{+} \doteq \vz + \delta \sigma \nabla \log p_{\sigma}(\vz) = \vz - \delta \sigma(\vSigma + \sigma^{2}\vI)^{-1}(\vz - \vu).
    \end{equation}
    最后，我们将自己限制在一个值得注意的特例中。设 \(\vu = \vzero\)，设 \(r \in [d]\) 是一个正整数，并设 \(\vSigma = \vU\vU^{\top}\) 对于某个 \(\vU \in \O(d, r)\)。定义 \(\cF_{\sigma, r}\) 为
    \begin{equation}
        \cF_{\sigma, r} = \{\vx \mapsto \hat{\vSigma}(\hat{\vSigma} + \sigma^{2}\vI)^{-1}\vx \colon \hat{\vSigma} = \hat{\vU}\hat{\vU}^{\top}, \hat{\vU} \in \O(d, r)\},
    \end{equation}
    即对应于具有真实均值 \(\vzero\) 和估计的真实协方差 \(\hat{\vU}\hat{\vU}^{\top} + \sigma^{2}\vI\) 的高斯分布的去噪器集合，这与真实模型通过支持估计（即 \(\hat{\vU}\) 与 \(\vU\)）不同。然后，如 \cite{wang2024diffusion} 所示，去噪问题
    \begin{equation}\label{eq:denoise}
        \argmin_{f_{\sigma} \in \cF_{\sigma, r}}\Ex\norm{\vz_{o} - f_{\sigma}(\vz_{\sigma})}_{2}^{2}
    \end{equation}
    等价于 PCA，即识别支持数据分布的 \(\hat{\vU} \in \O(d, r)\)。通过这种方式，通过去噪进行分布学习可以被视为非线性数据分布的 PCA 的推广。

\begin{exercise}
请证明问题 \eqref{eq:denoise} 等价于概率主成分分析（PCA）问题。以下是一些提示： \\
(i) 矩阵求逆引理：$(\vI + \vQ\vQ^\top)^{-1} = \vI - \vQ(\vI + \vQ^\top\vQ)^{-1}\vQ^\top$ 对于任何 $\vQ \in \R^{d\times r}$。 \\
(ii) 对于任何 $\vU \in \O(d, r)$，有 $\mathbb{E}_{\vg \sim \mathcal{N}(\vzero, \vI)}[\|\vU^\top \vg\|^2] = \sum_{i=1}^r \mathbb{E}_{\vg}[\|\vu_i^\top\vg\|^2] = r$，其中 $\vu_i$ 表示 $\vU$ 的第 $i$ 列。 \\
(iii) 你可以条件化 $\vz_o$。
\end{exercise}

\subsection{去噪低维高斯混合}

% \yima{%将这作为一个单独的小节，以便与第3.3.3节配对。%考虑改进这一小节的内容：
% \begin{enumerate}
%     %\item 陈述关于此类模型的采样复杂度的最佳已知结果，从记忆到泛化等。与第3.3.3节建立联系。 \DP{已知对于\textit{各向同性}高斯分布的采样复杂度，但不是低秩高斯分布（Peng的论文也使用了一些阈值操作符，而不是真正的去噪器，但有一些更进一步的工作，例如Kulin Shah和Sitan Chen的工作）。同样，记忆-泛化转换也不是已知的，尽管我们正在研究它。}
%     %\item 使用此类模型说明之前去噪和采样算法的具体应用。 \DP{在前面的单高斯简单情况下已经完成。}
%     \item 添加一个特定混合分布的模拟示例，例如与第3.3.3节中使用的相同模型：在 $\mathbb{R}^3$ 中的一个平面和两条直线。同时设计并留下一些编程练习，让学生实现用于学习类似简单分布的扩散过程。
% \end{enumerate}
% }

%\begin{example}\label{example:gmm_denoising}
    在本节中，我们将通过一个非常重要的分布族来分析去噪过程，其中数据来自低秩高斯混合模型，估计器试图估计支撑子空间。图 \ref{fig:denoising-twolines-oneplane} 已经数值地展示了这种分布的一个特定示例的去噪过程。我们在前一章 \ref{ch:linear-independent} 中研究了这个分布族的一些特殊情况，并将在下一章 \ref{ch:representation} 中看到，这个分布族可以用来很好地建模大多数实际重要的通用分布。


更具体地说，假设数据是从以下分布中抽取的
    \begin{equation}
        \vz_{o} \sim \frac{1}{K}\sum_{k = 1}^{K}\dNorm(\vzero, \vU_{k}\vU_{k}^{\top}),
    \end{equation}
    其中 \(\vU_{[K]} = (\vU_{k})_{k = 1}^{K}\) 是 \(d \times p\) 的正交矩阵。
    通过使用 Tweedie 公式（留作练习）进行计算，可以得到 \(\vz_{o}\) 在噪声水平 \(\sigma\) 下的最优去噪器为
    \begin{equation}
        f_{\sigma}^{\star}(\vz) = \frac{\sum_{k = 1}^{K}\phi(\vz \mid \vzero, \vU_{k}\vU_{k}^{\top} + \sigma^{2}\vI)\cdot (\vU_{k}\vU_{k}^{\top} + \sigma^{2}\vI)^{-1}\vz}{\sum_{k = 1}^{K}\phi(\vz \mid \vzero, \vU_{k}\vU_{k}^{\top} + \sigma^{2}\vI)},
    \end{equation}
    其中 \(\phi(\vx \mid \vmu, \vSigma)\) 是均值为 \(\vmu\) 和 \(\vSigma\) 的正态分布的密度函数在点 \(\vx\) 处的值。
    现在，假设我们使用一个与具有 \(\wh{K}\) 个分量的类似测度对应的去噪器类，具有估计的子空间 \(\hat{\vU}_{k}\)，即属于类 \(\cF_{\wh{K}, \sigma}\)，包含以下函数
    \begin{equation}\label{eq:estimated_lowrank_denoiser}
        f_{\sigma}(\vz \mid \hat{\vU}_{[\hat{K}]}) \doteq \frac{\sum_{k = 1}^{\hat{K}}\phi(\vz \mid \vzero, \hat{\vU}_{k}\hat{\vU}_{k}^{\top} + \sigma^{2}\vI)\cdot (\hat{\vU}_{k}\hat{\vU}_{k}^{\top} + \sigma^{2}\vI)^{-1}\vz}{\sum_{k = 1}^{K}\phi(\vz \mid \vzero, \hat{\vU}_{k}\hat{\vU}_{k}^{\top} + \sigma^{2}\vI)}.
    \end{equation}
    最近的工作，如 \cite{wang2024diffusion} 所示，通过在去噪目标上进行梯度下降，可以学习 \(\hat{\vU}_{[\hat{K}]}\)，并且在有足够的样本 \(N\) 和 \(K = \hat{K}\)（假设）的情况下，我们有 \(\hat{\vU}_{k}\hat{\vU}_{k}^{\top} \approx \vU_{k}\vU_{k}^{\top}\)。然而，这还不是全部——事实上，这个去噪器的表达式，\eqref{eq:estimated_lowrank_denoiser}，实际上几乎是以神经网络层的形式出现的。为了看到这一点，注意我们可以计算（证明留作练习）：
    \begin{equation}
        f_{\sigma}(\vz \mid \hat{\vU}_{[\hat{K}]}) = \sigma^{-2}\vz - \frac{\sigma^{-2}}{1 + \sigma^{2}}\cdot\frac{\sum_{k = 1}^{\hat{K}}\phi(\vz \mid \vzero, \hat{\vU}_{k}\hat{\vU}_{k}^{\top} + \sigma^{2}\vI)\hat{\vU}_{k}\hat{\vU}_{k}^{\top}\vz}{\sum_{k = 1}^{K}\phi(\hat{\vz} \mid \vzero, \hat{\vU}_{k}\hat{\vU}_{k}^{\top} + \sigma^{2}\vI)}.
    \end{equation}
    这类似于变换器中的注意力层的缩放，但有以下几点不同：
    \begin{itemize}
        \item 只有一个“标记”，尽管这更多是数据模型的后果。
        \item 每个 \(K\) 个头部的查询-键-值矩阵被简单的线性投影矩阵 \(\vU_{k}\) 替换。
    \end{itemize}
    这不是巧合，也不是偶然！在接下来的第 \ref{sec:lossy_compression} 节和第 \ref{ch:representation} 章中，我们将从压缩的角度给出为什么类似变换器的网络会使用这些去噪器的原因。
%\end{example}

让我们花一点时间反思一下我们所取得的成就：
\begin{itemize}
    \item 我们发现了一种通过参数化函数类中的多个去噪器来学习分布编码的方法。
    \item 我们发现了一种使用该参数化编码（例如，去噪器）来迭代地减少高熵噪声样本的熵，直到它落在数据分布上的方法。
\end{itemize}
这两种方法在某种意义上实际上满足了在 \eqref{sub:min_entropy} 中提出的候选方法。然而，所有这些都或多或少是隐式的：去噪器仅形成分布的隐式编码，采样仅在分布级别上进行。在下一节中，我们将讨论一个使我们的分布模型更明确和计算上可访问的框架。

\begin{exercise}
    完成本小节中提到的所有证明。具体来说：
    \begin{enumerate}
        \item 通过Tweedie公式（定理 \ref{thm:tweedie}）证明，对应于高斯混合分布 \(\mu = \frac{1}{K}\sum_{k = 1}^{K}\dNorm(\vzero, \vU_{k}\vU_{k}^{\top})\) 的最优去噪器，其中 \(\vU_{k} \in \O(d, p)\)，是
        \begin{equation}
            f_{\sigma}^{\star}(\vx \mid \vU_{[K]}) = \frac{\sum_{k = 1}^{K}\phi(\vx \mid \vzero, \vU_{k}\vU_{k}^{\top} + \sigma^{2}\vI)(\vU_{k}\vU_{k}^{\top} + \sigma^{2}\vI)^{-1}\vx}{\sum_{i = 1}^{K}\phi(\vx \mid \vzero, \vU_{k}\vU_{k}^{\top} + \sigma^{2}\vI)}
        \end{equation}
        其中 \(\phi(\vx \mid \vmu, \vSigma)\) 是 \(\dNorm(\vmu, \vSigma)\) 在 \(\vx\) 处的密度。
        \item 证明Woodbury矩阵恒等式，即
        \begin{equation}
            (\vA + \vU\vC\vV)^{-1} = \vA^{-1} - \vA^{-1}\vU(\vC^{-1} + \vV\vA^{-1}\vU)^{-1}\vV\vA^{-1}
        \end{equation}
        其中 \(\vA \in \R^{n \times n}\) 是可逆矩阵，\(\vU \in \R^{n \times k}\)，\(\vV \in \R^{k \times n}\)，且 \(\vC \in \R^{k \times k}\) 是可逆矩阵。
        \item 使用Woodbury矩阵恒等式证明
        \begin{align}
            (\vU_{k}\vU_{k}^{\top} + \sigma^{2}\vI)^{-1} 
            &= \sigma^{-2}\vI - \sigma^{-4}\vU_{k}(\vI + \sigma^{-2}\vI)^{-1}\vU_{k}^{\top} \\ 
            &= \sigma^{-2}\vI - \frac{\sigma^{-4}}{1 + \sigma^{-2}}\vU_{k}\vU_{k}^{\top} \\ 
            &= \sigma^{-2}\vI - \frac{\sigma^{-2}}{1 + \sigma^{2}}\vU_{k}\vU_{k}^{\top} \\ 
            &= \sigma^{-2}\bp{\vI - \frac{1}{1 + \sigma^{2}}\vU_{k}\vU_{k}^{\top}}.
        \end{align}
        %\DP{不需要包含简化过程的所有步骤，这只是我推导出来的方式。}
        \item 使用上述逆矩阵的形式，将最优去噪器写为
        \begin{equation}
            \scalebox{0.95}{\(\displaystyle f_{\sigma}^{\star}(\vx \mid \vU_{[K]}) = \sigma^{-2}\vx - \frac{\sigma^{-2}}{1 + \sigma^{2}}\cdot\frac{\sum_{k = 1}^{K}\phi(\vx \mid \vzero, \vU_{k}\vU_{k}^{\top} + \sigma^{2}\vI)\vU_{k}\vU_{k}^{\top}\vx}{\sum_{k = 1}^{K}\phi(\vx \mid \vzero, \vU_{k}\vU_{k}^{\top} + \sigma^{2}\vI)}.\)}
        \end{equation}
    \end{enumerate}
\end{exercise}

\begin{exercise}\label{exercise:vp_tweedie_denoising}
    考虑新的数据生成过程
    \begin{equation}
        \vz_{\sigma} = \alpha(\sigma)\vz_{o} + \sigma \vg, \qquad \vg \sim \dNorm(\vzero, \vI),
    \end{equation}
    其中 \(\alpha \colon \R_{> 0} \to \R_{> 0}\) 是一个可积的正的非增函数。找到Tweedie公式的适当版本，例如计算 \(\Ex[\vz_{o} \mid \vz_{\sigma}]\)，并计算该过程的增量去噪迭代，例如计算 \(\Ex[\vz_{\sigma - \delta} \mid \vz_{\sigma}]\)。我们指出，这不是我们框架的任意扩展：它实际上扩展了框架以涵盖所有流行的扩散模型设置，包括那些使用所谓的“方差保持”和Ornstein-Uhlenbeck扩散（要求 \(\sigma \in [0, 1)\) 并取 \(\alpha(\sigma) = \sqrt{1 - \sigma^{2}}\)）或“流匹配”过程（要求 \(\sigma \in [0, 1]\) 并取 \(\alpha(\sigma) = 1 - \sigma\)）。
\end{exercise}

% 但是，从这个解释中我们可以观察到，该过程容易陷入密度的局部模式或局部最大值。为了避免这种情况，通常的采样算法在每次迭代中添加适量的高斯噪声，得到迭代
% \begin{equation}
%     \vz^{(m + 1)} = \vz^{(m)} + \delta \sigma^{(m)}\nabla \log p_{\sigma^{(m)}}(\vz^{(m)}) + \sqrt{2\delta \sigma^{(m)}}\vg^{(m)}
% \end{equation}
% 其中 \(\vg^{(m)}\) 是独立同分布的标准高斯变量。

% \yima{再次提醒，尽量避免对事物进行评论。我们不是在写论文，而是在写一本书。直接描述接下来可以或应该做什么。留待以后再进行引用或评论。书籍应该基本上自成一体。}
% 学习 \(\vs_{\sigma} = \nabla \log p_{\sigma}\) 是有用的，因为它提供了一种高效的学习参数化统计模型的方法 \DP{引用：aapo}，同时也提供了从扰动分布 \(p^{\sigma}\) 中采样的方法（如果 \(\sigma\) 较小或接近于 \(0\)，则也可以从真实分布 \(p\) 中采样）。 一种经典的采样方法是通过所谓的 \textit{Langevin 动力学} 获得的，最近被更现代的 \textit{扩散模型} 所取代。在接下来的内容中，我们将简要概述这两种方法。

% 使用 Langevin 动力学从分布 \(p_{\sigma}\) 中采样 \DP{不确定这里正确的引用是什么... 最坏的情况，查看 Lee/Risteski 博客中的来源}，我们使用以下迭代过程：
% \begin{equation}\label{eq:langevin}
%     \vz_{t + 1} = \vz_{t} + \eta_{t} \vs_{\sigma}(\vz_{t}) + \sqrt{2\eta_{t}} \vg_{t}, \qquad \vg_{t} \sim \dNorm(\vzero, \vI), \quad \forall t \in [T].
% \end{equation}
% 这里 \(\eta_{t}\) 是步长，\(\vg_{t}\) 是独立于整个轨迹直到时间步 \(t\) 的采样。 这个迭代过程的一些简要直觉如下。 当从期望分布 \(p_{\sigma}\) 中采样时，我们希望将我们的迭代 \(\vz_{t}\) 移向 \(p_{\sigma}\) 的高概率区域（等价于 \(\log p_{\sigma}\)，我们使用其梯度，因为它不需要计算 \(p_{\sigma}\) 本身中存在的难以处理的归一化常数）。 因此，我们在 \(\log p_{\sigma}\) 上采取梯度步。 但我们不想仅仅将样本集中在 \(\log p_{\sigma}\) 的局部或全局最优值上，因此为了将当前迭代从这些最优值中扰动出来，我们添加一些适当缩放的噪声。 可以证明，迭代 \eqref{eq:langevin} 将收敛到一个具有目标分布的随机变量。

% \begin{theorem}
% \yima{将上述 Langevin 动力学收敛到期望的（Gibbs）分布作为定理呈现。}    
% \end{theorem}

%\DP{图示说明低维密度的平滑形式上的梯度下降：类似于 MAE 论文中的图片，但数据是弯曲的} \DP{一系列图片：通用数据、线性数据、分段线性数据、将数据转换为线性结构}

% \paragraph{一个迭代去噪的过程。}
% \yima{在介绍近似和采样数据分布的解决方案之前，请描述该方法并为它设置舞台。描述从标准高斯分布到数据分布的迭代去噪过程。精确描述序列 \(\vs_{\sigma^{\ell}}(\vz_{t}^{\ell})\) 是什么以及如何学习它们，通过应用先前的结果。}

% \yima{去噪过程何时停止？为什么不收敛到样本的经验分布，这似乎具有最低的内在维度，即 0？}

% \subsection{退火Langevin动力学}
% 我们在这里对Langevin动力学的主要问题是它不允许我们从现实世界的数据分布中采样，即那些高维但具有低维结构的分布。特别是，这些分布没有明确定义的密度，因此Langevin动力学也不明确定义。出于这个原因，\DP{引用：song \& ermon 2019}提出了退火Langevin动力学的一种形式：运行\eqref{eq:langevin}的轨迹，使用越来越小的\(\sigma = \sigma^{\ell}\)，将前一个轨迹的终端状态作为当前轨迹的输入状态，直到输出分布接近数据分布，即
% \begin{align}
%     \vz_{t + 1}^{\ell} 
%     &= \vz_{t}^{\ell} + \eta_{t}^{\ell} \vs_{\sigma^{\ell}}(\vz_{t}^{\ell}) + \sqrt{2\eta_{t}^{\ell}} \vg_{t}^{\ell}, \qquad \vg_{t}^{\ell} \sim \dNorm(\vzero, \vI), \\ 
%     \vz_{1}^{\ell}
%     &= \vz_{T + 1}^{\ell - 1}, \qquad \forall t \in [T], \quad \forall \ell \in [L].
% \end{align}
% 通过适当选择\(\sigma^{\ell}\)，我们得到以下简化过程：
% \begin{equation}\label{eq:diffusion_sampling}
%     \vz^{\ell + 1} = \vz^{\ell} + \eta^{\ell} \vs_{\sigma^{\ell}}(\vz^{\ell}) + \sqrt{2\eta^{\ell}} \vg^{\ell}, \qquad \vg^{\ell} \sim \dNorm(\vzero, \vI), \quad \forall \ell \in [L].
% \end{equation}
% 这是一个与\textit{扩散模型}相关的迭代方案，假设得分\(\nabla \log p^{\sigma}\)估计得当，即使分布是低维的，也能证明从数据分布中采样。 \yima{这个陈述的任何参考文献？}

% 特别地，迭代\eqref{eq:diffusion_sampling}提供了一种从任意分布中\textit{去噪采样}的方法。也就是说，通过使用\eqref{eq:tweedie_simple}，我们可以将\eqref{eq:diffusion_sampling}写为
% \begin{align}
%     \vz^{\ell + 1} 
%     &= \vz^{\ell} + \eta^{\ell}\left(\frac{\vmu_{\sigma^{\ell}}(\vz^{\ell}) - \vz^{\ell}}{(\sigma^{\ell})^{2}}\right) + \sqrt{2\eta^{\ell}}\vg^{\ell}, \quad \vg^{\ell} \sim \dNorm(\vzero, \vI), \quad \forall \ell \in [L], \\
%     &= \left(1 - \frac{\eta^{\ell}}{(\sigma^{\ell})^{2}}\right)\vz^{\ell} + \frac{\eta^{\ell}}{(\sigma^{\ell})^{2}}\vmu_{\sigma^{\ell}}(\vz^{\ell}) + \sqrt{2\eta^{\ell}}\vg^{\ell}.
% \end{align}
% \yima{可能需要呈现Jun Zhu小组的结果，说明如何使所有量明确，例如方差等。}
% 特别地，每次迭代都是当前迭代和最优去噪器的凸组合，逐渐向去噪样本移动，同时添加一些适当缩放的噪声以沿后验分布采样。因此，扩散模型通过估计最优去噪器来学习分布，并通过迭代去噪进行采样。

% \yima{将过程描述为可实现算法的伪代码。} \DP{}

% % \DP{大纲和一些旧笔记}这些也是在高维空间中学习具有非常低的内在维度分布的关键方法。更具体地说，去噪（和得分函数）、完成或预测等。

% % 经验贝叶斯去噪、Tweedie公式、得分函数、Langevin动力学、MCMC等。（可以参数化得分函数，但不需要...去噪在一般情况下是明确定义的。 \yima{不理解最后一句话。}）

% \DP{TODO: DP}

% \subsection{低秩高斯的去噪}
% \yima{请注意，上述框架适用于任意分布。一般来说，我们不知道得分函数\(\boldsymbol{s}_\sigma\)的确切形式。当然，可以使用神经网络近似它，就像在实践中经常做的那样。这里我们展示，当底层分布可以用低秩高斯建模时，得分函数具有显式形式，这有助于我们理解现代深度网络的结构。}

% \yima{为什么去噪过程不会收敛到（离散）经验分布，其内在维度最低，为0？如果没有严格的证明，至少提供一些合理的解释。可能引导向损失压缩方法。}

% 扩散模型和去噪器被认为学习了分布。一个自然的问题是，它们所学习的信息在某种意义上是否等同于其他经典方法学习分布时所学习的正确信息。事实证明，在某些受限的情况下，可以证明答案是肯定的。 \DP{继续...}

% \DP{大纲}
% 从扩散和去噪的角度解释PCA，针对单个低秩高斯分布。

% 低秩高斯混合的扩散和去噪，以及子空间聚类。

% \DP{TODO: DP}

\section{有损压缩} \label{sec:lossy_compression}

%\yima{有损编码不仅仅是量化。它允许我们通过简单地使用经验分布作为最小化熵的最优解来避免病态解。如何详细阐述这一点？}

敏锐的读者可能会注意到似乎存在一个困难：如果我们最终从有限样本中学习分布，并且我们的去噪器函数类包含足够多的函数，我们如何确保我们从\textit{真实}分布中采样，而不是从其他可能以高概率产生这些有限样本的分布中采样？特别是，从几何上讲，生成的有限样本集的经验分布是{\em 最小体积}的数据集，支持该有限样本集。\footnote{一组离散样本都是零维的。} 例如，图 \ref{fig:1d-line} 中所示的两组采样数据点在几何上是相同的：每组包含八个点，每个点出现的频率相同。唯一的区别是，对于第二个数据集，某些点足够“接近”，可以被视为在其各自的“簇”周围具有更高的密度。

\begin{figure}
    \centering
    \includegraphics[width=0.7\linewidth]{figures/one-dim-distribution.png}
    \caption{一条线上的八个点。}
    \label{fig:1d-line}
\end{figure}

如果（几何）体积不是衡量（有限采样经验）分布的简约性或压缩性的合适标准，那么也许与这些样本 $p(\vx_i)$ 的真实概率（密度或分布）相关的{熵}，或{\em 最小编码率}，是一个更好的衡量标准。尽管许多结果，如定理 \ref{thm:distribution-convergence}，确保了在渐近情况下收敛到真实分布（因此熵），但在非渐近情况下，当我们仅给定一组固定样本时，如何同时推断分布并正确地进行压缩，这一点并不明确，如图 \ref{fig:1d-line} 所示的情况。实证观察表明，采样算法 \ref{alg:iterative_denoising}（使用扩散/去噪模型）倾向于超越仅从训练集生成样本。目前，这种现象发生的原因和时机仍然是一个开放问题。

此外，除了数据分布的熵之外，实际上还有其他与实现任何给定数据集的编码-解码方案相关的根本约束和成本。本节尝试从构建的角度解决这些问题，探讨从有限样本集学习和表示数据分布的问题。

\subsection{有损编码率和熵}
%\yima{从可计算（熵）到可实现。也就是说，除了可计算或可实现的编码方案外，还应该存在可实现的解码方案。对于实值数据，有损量化变得必要。}

给定一个采样数据集 $\X = [\vx_1, \ldots, \vx_N]$，如何设计一个可以在有限内存和计算资源的机器上实现的压缩方案？请注意，即使表示一个一般的实数也需要无限位或无限个比特。
\begin{example}
    考虑一个离散分布 $\X = [e, \pi]$，其中每个值 $e$ 或 $\pi$ 的概率均为 $1/2$。该分布的熵为 $H =1$。也就是说，可以分别用一个比特的数字 $0$ 或 $1$ 来编码这两个数。但是，你能在有限状态机上实现这个代码的解码方案吗？答案实际上是否定的，因为精确描述任何一个数都需要无限多个比特。
\end{example}

因此，通常不可能有一个编码和解码方案能够精确地再现任意实值分布的样本。\footnote{也就是说，如果想要精确地编码这样的样本，唯一的方法是记住每一个样本。} 但是，如果不能解码来自同一分布的样本，那么编码一个分布就几乎没有实际价值。

因此，为了确保压缩编码/解码方案是可计算的，并且可以用有限的内存和计算资源实现，我们需要量化样本 $\z$ 并仅编码到一定的精度，例如 $\epsilon > 0$。{\em 通过这样做，本质上，我们将任何两个距离小于 $\epsilon$ 的数据点视为等价。} 更精确地说，我们希望考虑编码方案
\begin{equation}
    \z \mapsto \hat \z
\end{equation}
使得量化引起的期望误差小于等于 $\epsilon$：$$\mathbb{E} [d(\z, \hat \z)] \le \epsilon.$$ 通常，距离选择为欧几里得距离或 2-范数。

当然，在所有满足上述约束的编码方案中，我们希望选择使结果编码率最小的那个。对于给定的随机变量 $\z$ 和精度 $\epsilon$，这个率被称为{\em 率失真}，记为 $R_{\epsilon}(\z)$。事实证明，对于一个实值连续随机向量，这个率是 $\z$ 的熵 $H(\z)$ 的一个可实现的近似，具体来说：
\begin{equation}
    R_{\epsilon}(\z) = \min_{p(\z, \hat \z): \mathbb{E} [d(\z, \hat \z)] \le \epsilon} H(\z) - H(\z \mid \hat{\z}),
\end{equation}
其中最小化是在所有满足失真约束 $\mathbb{E}[d(\z, \hat \z)] \le \epsilon$ 的（联合）分布 $p(\z, \hat{\z})$ 上进行的。从定义中，我们知道 $R(\z \mid \epsilon)$ 是一个关于 $\epsilon$ 的{\em 非增}函数。

\begin{remark}
    注意，在信息论中，量 $H(\z) - H(\z\mid \hat \z)$ 被称为两个变量 $\z$ 和 $\hat \z$ 之间的{\em 互信息} $I(\z; \hat \z)$。因此，最小编码率是在所有满足约束 $\mathbb{E} [d(\z, \hat \z)] \le \epsilon$ 的分布中，使 $\z$ 和 $\hat \z$ 之间的互信息最小化时实现的。更多细节请参见 \cite{Cover-Thomas}。
\end{remark}

\begin{remark}
    注意到给定一组数据点 $\X = [\x_1,\ldots, \x_N]$，我们总可以将它们解释为从一个均匀离散分布中抽取的样本，该分布在这 $N$ 个向量上的概率相等，即 $1/N$。这种分布的熵为 $H(\X) = \frac{1}{N}\log_2 N$。\footnote{再次注意，即使我们可以用这种编码率对这些向量进行编码，我们也不能以任意精度对其进行解码。} 然而，即使 $\X$ 是其样本上的均匀分布，如果这些样本分布不均匀且许多样本紧密聚集在一起，那么使用有损编码方案可以实现的编码率 $R_{\epsilon}(\X)$ 可能远低于 $H(\X)$。\footnote{尽管如此，可以证明，对于这种离散均匀分布，当 $\epsilon$ 足够小时，我们总是有 $H(\X) = R(\X \mid \epsilon)$。}
\end{remark}

\begin{example}
有时，我们可能会遇到相反的情况，即我们首先固定编码率，然后尝试找到一个最小化失真的编码方案。例如，我们只想使用固定数量的代码对从一个正方形中抽取的点进行编码，并且我们想知道如何设计这些代码，使得在编码/解码方案中失真最小。例如，给定一个单位正方形上的均匀分布，我们想知道如何用 $n$ 位来编码从该分布中抽取的点，使得编码的精度最高。很容易看出，这个问题等价于如何将 $2^n$ 个最小半径的圆盘打包以覆盖这个单位正方形。图 \ref{fig:seven-circles-packing} 显示了分别用 7、15 和 24 个圆盘覆盖正方形的最优打包方式。
\begin{figure}
    \centering
    \includegraphics[width=0.8\linewidth]{figures/Optimal-solutions-for-7-15-and-24-circle-packing-in-a-square.png}
    \caption{7、15 和 24 个圆盘覆盖正方形的最优解。}
    \label{fig:seven-circles-packing}
\end{figure}
\end{example}

对于一般的分布，通常不可能找到其率失真函数的解析形式。人们通常求助于数值计算\footnote{感兴趣的读者可以参考 \cite{Blahut-1972} 中的一个经典算法，该算法可以数值计算离散分布的率失真函数。}。然而，正如我们将看到的，在我们的上下文中，我们通常需要知道率失真作为一组数据点或其表示的显式函数。这是因为我们希望将编码率作为表示质量的度量。显式的解析形式使得确定如何转换数据分布以改进表示变得容易。因此，我们应该使用那些率失真函数具有显式解析形式的分布。为此，我们从最简单且最重要的分布族开始。

\subsection{有损编码率对于低维高斯分布}
%\yima{有损编码长度，通过球填充进行编码和解码，高斯分布或低维子空间的率失真。}
现在考虑我们给定了一组来自任意分布的数据样本 $\X = \{\x_1, \ldots, \x_N\}$，\footnote{或者这些数据点可以被视为一个（经验）分布本身。} 我们希望提出一种构造性方案，可以将数据编码到一定的精度，例如
\begin{equation}
    \x_i \mapsto \hat \x_i, \quad \mbox{subject to} \quad \|\x_i - \hat \x_i\|_2 \le \epsilon.
\end{equation}
不失一般性，假设 $\X$ 的均值为零 $ \sum_{i} \x_i/N = \boldsymbol{0}$。在没有任何关于 $\X$ 背后的分布性质的先验知识的情况下，我们可以将 $\X$ 视为从高斯分布 $\mathcal{N}(\boldsymbol{0}, {\boldsymbol{\Sigma}})$ 中采样，其协方差为\footnote{已知给定固定方差，高斯分布达到最大熵。也就是说，它给出了在可能的编码率方面最坏情况的上限。}：
\begin{equation}
{\boldsymbol{\Sigma}} = \frac{1}{N} \X\X^\top. 
\end{equation}
注意，几何上 ${\boldsymbol{\Sigma}}$ 表征了一个椭球区域，其中大部分样本 $\x_i$ 位于该区域。

我们可以将 $\hat \X = [\hat \x_1,\ldots, \hat \x_N]$ 视为 $\X = [\x_1, \ldots, \x_N]$ 的噪声版本：
\begin{equation}
    \hat \x_i = \x_i + \boldsymbol{n}_i
\end{equation}
其中 $\boldsymbol{n}_i$ 被近似为独立的高斯噪声 $\boldsymbol{n}_i  \sim \mathcal{N}(\boldsymbol{0}, {\epsilon^2}  \boldsymbol{I}/{D})$。那么 $\hat \x_i$ 的协方差由下式给出：
\begin{equation}
\hat{\boldsymbol{\Sigma}} = \mathbb{E}\left[\frac{1}{N} \hat \x_i \hat \x_i^\top\right] = \frac{\epsilon^2}{D} \boldsymbol{I} + \frac{1}{N} \X\X^\top.
\end{equation}
注意，由向量 $\x_i$ 张成的区域的体积与协方差矩阵的行列式的平方根成正比：
\begin{equation}
    \mbox{体积}(\hat \x_i) \propto \sqrt{\det \big(\hat{\boldsymbol{\Sigma}}\big)} = \sqrt{\det\left(\frac{\epsilon^2}{D} \boldsymbol{I} + \frac{1}{N} \X\X^\top \right)}.
\end{equation}
由每个随机向量 $\boldsymbol{n}_i$ 张成的区域的体积与下式成正比：
\begin{equation}
  \mbox{体积}(\boldsymbol{n}_i) \propto   \sqrt{\det\left(\frac{\epsilon^2}{D} \boldsymbol{I} \right)}.
\end{equation}

\begin{figure}
    \centering
    \includegraphics[width=\linewidth]{figures/Gaussian-compression.png}
    \caption{在数据向量张成的区域中填充 $\epsilon$-球。空间的体积越大，需要的球越多，因此需要更多的比特来编码和枚举这些球。每个实值向量 $\z$ 可以编码为它落入的球的编号。}
    \label{fig:ball-packing}
\end{figure}

为了编码落入 $\hat \x_i$ 张成的区域的向量，我们可以用不重叠的半径为 $\epsilon$ 的球覆盖该区域，如图 \ref{fig:ball-packing} 所示。当由 $\hat \x_i$ 张成的区域的体积显著大于 $\epsilon$-球的体积时，我们需要覆盖该区域的球的总数大约等于两个体积的比值：
\begin{equation}
    \# \,\epsilon\mbox{-球} = \frac{\mbox{体积}(\hat \x_i)}{\mbox{体积}(\boldsymbol{n}_i)} = \sqrt{\det\left(\boldsymbol{I} + \frac{D}{N\epsilon^2} \X\X^\top  \right)}.
\end{equation}

如果我们使用二进制数来标记感兴趣区域内的所有 $\epsilon$-球，所需的总二进制位数为 %\DP{建议使用符号 \(R_{\eps}\)，因为“条件”符号在将来会过度使用。}
\begin{equation}
    R_{\epsilon}(\X) \doteq \log_2 (\# \,\epsilon\mbox{-球}) = \frac{1}{2} \log \det \left(\boldsymbol{I} + \frac{D}{N\epsilon^2} \X\X^\top \right).
    \label{eqn:rate-Gaussian}
\end{equation}

\begin{example}
图 \ref{fig:ball-packing} 显示了一个具有椭球支撑的 2D 分布示例——近似于 2D 高斯分布的支撑。该区域被大小为 $\epsilon$ 的小球覆盖。所有的小球从 1 到 n 编号。然后，对于该区域中的任何向量 $\z$，我们只需要确定它最接近哪个 $\epsilon$-球心，记为 ball$(\z)$。为了记住 $\z$，我们只需要记住该球的编号，这需要 $\log(n)$ 位来存储。如果我们需要从这个数字解码 $\z$，我们只需将 $\hat \z$ 取为该球的中心。这导致了一个显式的编码和解码方案：
\begin{equation}
    \z \longrightarrow \mbox{ball}(\z) \longrightarrow \hat \z = \mbox{ball}(\z) 的中心。
\end{equation}
人们可以将这些球心称为编码方案的代码簿或字典中的“代码”。很容易看出，这种（有损）编码-解码方案的精度约为球的半径 $\epsilon$。显然，$R_{\epsilon}(\Z) = \log(n)$ 是使用这种编码方案编码每个向量 $\z$ 的球编号所需的平均位数，因此与该方案相关的编码率。

\end{example}

从上述推导中，我们知道编码率 $R_{\epsilon}(\X)$ 可以通过显式的编码（和解码）方案实现。它有两个有趣的性质：
\begin{itemize}
    \item 首先，人们可能会注意到 $R_{\epsilon}(\X)$ 与高斯源的速率失真函数非常相似 \cite{Cover-Thomas}。事实上，当 $\epsilon$ 很小时，上述表达式是高斯源速率失真的一个近似值，正如 \cite{MaY2007-PAMI} 指出的那样。
    \item 其次，如果假设数据 $\X$ 来自一个线性子空间，也可以导出相同的编码率 $R_{\epsilon}(\X)$。这可以通过适当地量化 $\X = \boldsymbol{U} \boldsymbol{\Sigma}\boldsymbol{V}^\top$ 的奇异值分解（SVD）并为 $\boldsymbol{U}$ 跨越的子空间中的向量构造一个有损编码方案来证明 \cite{MaY2007-PAMI}。
\end{itemize}
在我们的上下文中，表达式 $R_{\epsilon}(\X)$ 非常基本：它是与从高斯分布或线性子空间中抽取的数据的显式和自然有损编码方案相关的编码率。正如我们将在下一章中看到的，这个公式在理解深度神经网络的架构中起着重要作用。

\subsection{聚类低维高斯混合}
如前所述，给定的数据集 $\X$ 通常具有低维的内在结构。因此，将其编码为一般的高斯分布会非常冗余。如果我们能够识别 $\X$ 中的这些内在结构，我们可以设计出更好的编码方案，给出更低的编码率。或者等价地，用于编码这种 $\X$ 的代码可以被压缩。我们将看到，压缩提供了一种统一的可计算方法来识别这些结构。在本节中，我们将通过最基础的低维结构族来展示这个重要思想：高斯混合或子空间。



\begin{example}
    图 \ref{fig:two-subspaces} 展示了一个例子，其中数据 $\X$ 分布在两个子空间（或低维高斯分布）周围。如果我们将它们视为一个单一的高斯分布并一起编码，那么与之相关的离散（有损）码本，由所有蓝色球体表示，显然是非常冗余的。我们可以尝试识别这两个子空间的位置，并设计一个仅覆盖这两个子空间的码本，即绿色球体。如果我们能够正确地将数据 $\X$ 中的样本划分到这两个子空间：$\X = \X_1 \cup \X_2$，其中 $\X_1 \in S_1$ 和 $\X_2 \in S_2$，那么数据的编码率将大大降低。这将提供一个更简洁、因此更理想的表示。

\end{example}

\begin{figure}
    \centering
    \includegraphics[width=0.5\linewidth]{figures/Two-subspaces.png}
    \caption{比较两种针对分布在两个子空间周围的有损编码方案。一种是为两个子空间张成的整个空间打包（蓝色）$\epsilon$-球体；另一种是在两个子空间周围的表格邻域内仅打包球体。后者显然具有更小的码本，并且对子空间上的样本产生更低的编码率。}
    \label{fig:two-subspaces}
\end{figure}

因此，更一般地说，如果数据来自任何子空间或低维高斯分布的混合，那么识别这些组件并根据这些组件的内在维度对数据进行编码将是可取的。事实证明，通过这种假设，我们不会失去太多的一般性，因为众所周知，高斯混合可以很好地逼近大多数一般分布 \cite{borkar2016gaussian}。

\paragraph{聚类问题。}
现在，对于这个特定的分布族，我们如何有效地从一组样本中识别出这些低维组件
\begin{equation}
    \X = \{\x_1, \x_2, \ldots, \x_N\},
\end{equation}
从中抽取？换句话说，给定整个数据集 $\X$，我们希望将其划分为多个，例如 $k$ 个子集：
\begin{equation}
    \X = \X_1 \cup \X_2 \cup \cdots \cup \X_k,
\end{equation}
其中每个子集仅包含从一个低维高斯分布或子空间中抽取的样本。请注意，根据具体情况，划分可以是确定性的或概率性的。如 \cite{ma2007segmentation} 所示，对于高斯混合，概率划分不会导致更低的编码率。因此，为了简单起见，我们这里仅考虑确定性划分。

\paragraph{通过有损压缩进行聚类。}
解决上述聚类问题的主要困难在于，我们通常不知道簇的数量 $k$，也不知道每个组件的维度。这个问题的研究历史悠久。教科书 \cite{GPCA} 对这个问题的不同方法进行了系统和全面的介绍。

为了找到解决这个问题的有效方法，我们首先需要理解和澄清为什么我们想要聚类。换句话说，与不聚类相比，聚类数据究竟带来了什么好处？我们如何衡量这种好处？从数据压缩的角度来看，正确的聚类应该导致更高效的编码（和解码）方案。

对于任何给定的数据集 $\X$，已经存在两种显而易见的编码方案作为基线。它们代表了编码数据的两种极端方式：
\begin{itemize}
\item 简单地将所有样本视为来自一个单一的高斯分布。相关的编码速率由下式给出：
\begin{equation}
    R_{\epsilon}(\X) = \frac{1}{2} \log \det \left(\boldsymbol{I} + \frac{D}{N\epsilon^2} \X\X^\top \right).
\end{equation}
\item 简单地分别记忆所有样本，为每个样本分配一个不同的数字。编码速率将是：
\begin{equation}
    R_0(\X) = \log(N). 
\end{equation}
\end{itemize}

请注意，任一编码方案都可以成为特定（极端）量化误差 $\epsilon$ 选择下的“最优”解：
\begin{enumerate}
    \item {\em 懒惰模式}：如果我们选择 $\epsilon$ 非常大，$\X$ 中的所有样本都可以被一个单一的球体覆盖。速率是  $\lim_{\epsilon \rightarrow \infty} R_\epsilon \rightarrow \frac{1}{2}\log\det (\boldsymbol{I}) = 0$。 
    \item {\em 记忆模式}：如果 $\epsilon$ 非常小，$\X$ 中的每个样本都被一个不同的 $\epsilon$-球体覆盖，因此总数是 $N$。速率是 $\lim_{\epsilon \rightarrow 0} R_\epsilon \rightarrow \log(N)$。
\end{enumerate} 
请注意，第一种方案对应于完全不关心分布的任何有趣特性的场景。我们不想为任何有意义的信息节省任何比特。我们称这种模式为“懒惰模式”。第二种方案对应于希望以极高的精度解码每个样本的场景。因此，我们最好“记忆”每个样本。我们称这种模式为“记忆模式”。 
\begin{figure}
    \centering
\includegraphics[width=0.9\linewidth]{figures/circle-packing.png}
    \caption{二维平面上的一组随机样本。考虑为每个样本分配一个以该样本为中心的 $\epsilon$-圆盘。样本的密度从左到右增加。}
    \label{fig:circle-packing}
\end{figure}
\begin{example}
为了了解何时偏好记忆模式，让我们考虑一组在二维平面上单位面积内随机分布的样本，假设 $N$ 个样本。\footnote{假设这些点是通过每单位面积 $N$ 个点的泊松过程绘制的。} 假设我们尝试设计一个具有固定量化误差 $\epsilon$ 的有损编码方案。这相当于在每个样本周围放置一个 $\epsilon$-圆盘，如图 \ref{fig:circle-packing} 所示。当 $N$ 较小时，所有圆盘相互重叠的概率为零。在这种情况下，需要且最优的码本大小为 $N$。当 $N$ 或密度达到某个临界值 $N_c$ 时，所有圆盘开始相互重叠并连接成一个覆盖整个平面的簇——这种现象被称为连续“渗流” \cite{Gilbert-1961,Mertens-Moore-2012}。当 $N$ 大于这个值时，圆盘重叠严重。圆盘的数量 $N$ 变得非常冗余，因为我们只想以给定的精度 $\epsilon$ 编码平面上的点。覆盖所有样本所需的圆盘数量远小于 $N$。\footnote{事实上，存在高效的算法来找到这样的覆盖 \cite{Booth-2001}。}
\end{example}

懒惰和记忆化策略在某种程度上是显而易见的，可能在理论或实际应用中兴趣不大。这两种方案在用于编码大量从具有{\em 紧凑且低维支持}的分布中抽取的样本时，都远非最优。事实证明，在这两种方案之间存在一个有趣的中间区域。

\begin{example}
\begin{figure}[t]
    \centering
\includegraphics[width=0.4\linewidth]{figures/Two-lines-and-plane.png} 
\includegraphics[width=0.9\linewidth]{figures/Coding-Rate.jpg}
    \caption{上图：从$\mathbb{R}^3$中的两条线和一个平面抽取的358个带噪声样本。下图：$\epsilon$变化对聚类结果和编码率的影响。红线标记了添加到样本中的高斯噪声的方差$\epsilon_0$。}
        \label{fig:two-lines-and-plane} 
    \label{fig:two-lines-and-plane-epsilon}
\end{figure}
图 \ref{fig:two-lines-and-plane} 显示了一个从$\mathbb{R}^3$中的两条线和一个平面抽取的带噪声样本的例子。我们注意到，从右侧的图 (c) 可以看出，随着$\epsilon$的增加，最优编码率单调递减，这与率失真函数的性质相符。图 (a) 和 (b) 显示了当$\epsilon$从非常小（接近零）变化到非常大（趋向无穷大）时，编码率最小时的最优聚类数。我们可以清楚地看到，在图的两端存在懒惰和记忆化策略。但在图 (b) 中，当量化误差$\epsilon$选择为接近真实噪声方差$\epsilon_0$的水平时，最优聚类数是代表两条线和一个子空间的“正确”数三。我们非正式地将这个中间区域称为“泛化区域”。请注意，在这些区域之间发生了明显的相变。\footnote{据我们所知，目前还没有严格的理论来解释这些相变行为。}   
\end{example}

从上述讨论和例子中，我们看到，当量化误差相对于样本密度\footnote{或样本密度相对于量化误差}处于适当范围内时，最小化有损编码率可以让我们发现样本数据的潜在（低维）分布。{\em 因此，量化，最初作为一种实用性选择，似乎成为从有限样本的经验分布中学习连续分布的必要手段。}尽管目前还没有严格的理论来解释这一现象，但在这里，为了学习目的，我们关心的是如何利用这一现象来设计算法，以找到正确的分布。

让我们使用图 \ref{fig:two-subspaces} 中显示的简单例子来说明基本思想。如果可以将所有样本$\X$划分为两个聚类$\X_1$和$\X_2$，分别包含$N_1$和$N_2$个样本，那么相关的编码率将是\footnote{我们在这里忽略了一些需要编码每个样本的隶属关系的开销位，例如通过霍夫曼编码。}
\begin{equation}
    R^c(\X\mid \boldsymbol{\Pi}) = \frac{N_1}{N}R_{\epsilon}(\X_1) + \frac{N_2}{N}R_{\epsilon}(\X_2),
\end{equation}
其中我们用$\boldsymbol{\Pi}$表示划分的隶属关系。如果划分尊重分布的低维结构，在这种情况下，$\X_1$和$\X_2$分别属于两个子空间，那么得到的编码率应该显著小于上述两种基本方案：
\begin{equation}
    R^c(\X \mid \boldsymbol{\Pi},\epsilon) \ll R_{\epsilon}(\X), \quad     R^c(\X \mid \boldsymbol{\Pi}, \epsilon) \ll R_0(\X). 
\end{equation}
一般来说，我们可以将聚类问题转化为一个最小化编码率的优化问题：
\begin{equation}
    \min_{\boldsymbol{\Pi}}   R^c(\X \mid \boldsymbol{\Pi}, \epsilon) = \sum_{i=1}^k \frac{N_i}{N}R_{\epsilon}(\X_i).
\end{equation}

\paragraph{聚类优化策略。}
剩下的问题是，我们如何优化上述编码率目标以找到最优的聚类。针对这一目标，有三种自然的方法：
\begin{enumerate}
    \item 我们可以从整个集合 $\X$ 作为一个单一的聚类（即懒惰策略）开始，然后搜索（例如随机地）将其划分，以便能够得到更小的编码率。
    \item 反之，我们可以从每个样本 $\x_i$ 作为自己的聚类（即记忆策略）开始，然后搜索合并那些能够导致更小编码率的聚类。
    \item 或者，如果我们能够将隶属度 $\boldsymbol{\Pi}$ 表示为某些连续参数，我们可以使用优化方法，如梯度下降（GD）。
\end{enumerate}
第一种方法在计算上不太吸引人，因为需要尝试的可能划分数量是样本数量的指数级。例如，将 $\X$ 划分为两个等大小子集的划分数量是 $N \choose N/2$，当 $N$ 变大时，这个数量会爆炸式增长。我们将在下一章 \ref{ch:representation} 中探讨第三种方法。在那里，我们将看到深度神经网络（特别是 Transformer）的作用如何与编码率目标相关联。

第二种方法最初由 \cite{ma2007segmentation} 提出。它展示了能够高效评估编码率（例如使用解析形式）的好处。有了它，数据的（低维）聚类可以通过最小化编码长度（MCL）的原则来找到，且效率和效果都相当高。

请注意，对于一个包含 $N_i$ 个样本的聚类 $\X_i$，编码 $\X_i$ 中所有样本所需的二进制位长度由下式给出：\footnote{实际上，编码长度的更准确估计是 $L(\X_i) = (N_i+D) R_\epsilon(\X_i)$，其中额外的位用于编码子空间的基 \cite{ma2007segmentation}。这里为了简化，我们省略了这一开销。}
\begin{equation}
    L(\X_i) = N_i R_\epsilon(\X_i). 
\end{equation}
如果我们有两个聚类 $\X_i$ 和 $\X_j$，如果我们希望将样本编码为两个独立的聚类，所需的二进制位长度为
\begin{equation*}
    L^c(\X_i, \X_j) = N_i R_\epsilon(\X_i) + N_j R_\epsilon(\X_j) - N_i \log\frac{N_i}{N_i+N_j} - N_j \log\frac{N_j}{N_i+N_j}. 
\end{equation*}
最后两项是根据霍夫曼编码编码样本隶属度所需的位数。

然后，对于任意两个独立的聚类 $\X_1$ 和 $\X_2$，我们可以根据两个编码长度之间的差异来决定是否合并它们：
\begin{equation}
    L(\X_i\cup \X_j) - L^c(\X_i, \X_j) 
\end{equation}
如果这个差值为负，这意味着如果我们将两个聚类合并为一个，编码长度会变得更小。这一简单事实导致了 \cite{ma2007segmentation} 提出的以下聚类算法：
\begin{algorithm}
    \caption{编码长度成对最速下降法}\label{alg:steepest_descent_coding_length}
    \begin{algorithmic}[1]
        \Require{\(N\) 数据点 \(\{\vx^{i}\}_{i = 1}^{N}\)}
        \Ensure{一组聚类 \(\cC\)}

\Procedure{配对最陡下降编码长度}{$\{\vx^{i}\}_{i = 1}^{N}$}
            \State{\(\cC \gets \{\{x^{i}\}\}_{i = 1}^{N}\)} \Comment{初始化 \(N\) 个簇 \(\vX_{i}\)，每个簇包含一个元素}
            \While{\(\abs{\cC} > 1\)}
                \If{\(\displaystyle \min_{\vX_{i}, \vX_{j} \in \cC}[L(\vX_{i} \cup \vX_{j}) -L^{c}(\vX_{i}, \vX_{j})] \geq 0\)} \Comment{如果任何合并都不会节省比特}
                    \State{\Return{\(\cC\)}} \Comment{提前返回 \(\cC\) 并退出}
                \Else
                    \State{\(\displaystyle \vX_{i^{\star}}, \vX_{j^{\star}} \gets \argmin_{\vX_{i}, \vX_{j} \in \cC}[L(\vX_{i} \cup \vX_{j}) - L^{c}(\vX_{i}, \vX_{j})]\)} \Comment{合并节省最多比特的簇}
                    \State{\(\displaystyle \cC \gets [\cC \setminus \{\vX_{i^{\star}}, \vX_{j^{\star}}\}] \cup \{\vX_{i^{\star}} \cup \vX_{j^{\star}}\}\)} \Comment{移除未合并的簇并添加合并后的簇}
                \EndIf
            \EndWhile
            \State{\Return{\(\cC\)}}  \Comment{如果所有合并都能节省比特，则返回一个簇}
        \EndProcedure
    \end{algorithmic}
\end{algorithm}

请注意，该算法是可行的，因为总共有 (成对) 比较和合并的次数约为 \(O(N^2\log N)\)。然而，由于其贪心性质，没有理论保证该过程会收敛到全局最优的聚类解。尽管如此，正如 \cite{ma2007segmentation} 中所报道的，实践中，这个看似简单的算法表现得非常好。图 \ref{fig:two-lines-and-plane} 中绘制的聚类结果实际上是由该算法计算得出的。

%\subsection{子空间聚类的应用}

\paragraph{示例：图像分割} 上述编码长度的度量和相关的聚类算法假设数据分布是 (低维) 高斯混合。尽管这看起来有些理想化，但在模型 (近似) 有效的情况下，该度量和算法仍然非常有用，甚至强大。

例如，自然图像通常由多个具有几乎均匀纹理的区域组成。如果我们从每个区域中取许多小窗口，它们应该类似于从 (低维) 高斯分布中抽取的样本，如图 \ref{fig:image-patch} 所示。图 \ref{fig:image-segmentation} 显示了将上述聚类算法直接应用于图像块进行图像分割的结果。有关将算法定制到图像分割问题的更多技术细节，请参见 \cite{Mobahi-IJCV2011}。
\begin{figure}
    \centering
    \includegraphics[width=0.5\linewidth]{figures/image-segmentation-tiles.png}
    \caption{大小为 \(w\times w\) 像素的图像块。}
    \label{fig:image-patch}
\end{figure}

\begin{figure}[t]
    \centering
    \includegraphics[width=0.95\linewidth]{figures/image-segmentation.png}
    \caption{基于将聚类算法应用于图像块的分割结果。}
    \label{fig:image-segmentation}
\end{figure}

%\subsection{通过最小增量编码长度进行分类}
%\href{http://people.eecs.berkeley.edu/~yima/psfile/MICL_SJIS.pdf}{监督分类} 的角度来看 (有损) 压缩。

\section{总结与注记}
{有损编码不仅限于量化。它允许我们通过简单地使用经验分布作为最小化熵的最优解来避免病态解。}

{从可实现的 (熵) 到可实现的编码方案。也就是说，除了可计算或可实现的编码方案外，还应该存在可实现的解码方案。对于实值数据，有损量化变得必要。}

\section{练习与扩展}

\begin{exercise}
请展示以下关于 $\log\det(\cdot)$ 函数的性质。

\begin{enumerate}
    \item 证明
    \begin{align*}
    f(\vX) = \log\det\left(\vX\right)  
    \end{align*}
    是一个凹函数。 （{\bf 提示：} 函数 $f(\vx)$ 是凸的当且仅当对于所有的 $\vx$ 和 $\vh$，函数 $f(\vx+t\vh)$ 也是凸的。）

\item 证明：
    \begin{align*}
        \log\det(\vI + \vX^\top\vX) = \log\det(\vI + \vX\vX^\top)
    \end{align*}

\item 设 $\vA \in \R^{n\times n}$ 是一个正定矩阵。请证明：
\begin{align}
\log\det\left(\vA \right) = \sum_{i=1}^n \log(\lambda_i),
\end{align}
其中 $\lambda_1,\lambda_2,\dots,\lambda_n$ 是 $\vA$ 的特征值。
\end{enumerate}
    
\end{exercise}

\begin{exercise}
    \sdb{关于证明高斯分布是最大熵的练习（在附录B中的一个备注中）}
\end{exercise}

\begin{exercise}
    \pw{在这里添加一个编程练习}
\end{exercise}

\end{document}